\documentclass[11pt,oneside]{article}

\usepackage[T1]{fontenc}
\usepackage[utf8]{inputenc}
\usepackage{lmodern}
\usepackage{microtype}
\usepackage{geometry}
\geometry{margin=1in,headheight=15pt}
\usepackage{amsmath,amssymb,amsthm,mathtools}
\usepackage{booktabs,longtable,tabularx,array}
\usepackage{enumitem}
\usepackage{xcolor}
\usepackage{xurl}
\usepackage{hyperref}
\usepackage{fancyhdr}
\usepackage{titlesec}
\usepackage{needspace}
\usepackage{listings}

\definecolor{proofgreen}{RGB}{15,105,68}
\definecolor{openblue}{RGB}{27,84,145}
\definecolor{classicalgray}{RGB}{85,92,99}
\definecolor{computeviolet}{RGB}{100,55,145}
\definecolor{codegray}{RGB}{248,248,248}
\definecolor{rulegray}{RGB}{92,101,111}

\hypersetup{
  colorlinks=true,
  linkcolor=openblue,
  citecolor=proofgreen,
  urlcolor=openblue,
  pdftitle={The Alaoglu--Erdos Integer-Exponent Conjecture: Polynomial Approximation, Integer-Valued Interpolation, and the Arithmetic Unit Barrier},
  pdfauthor={OpenAI, continuation of the ProveIt unified report},
  pdfsubject={Polynomial approximation methods for the Alaoglu--Erdos conjecture}
}
\urlstyle{same}

\pagestyle{fancy}
\fancyhf{}
\lhead{Alaoglu--Erd\H{o}s conjecture}
\rhead{Polynomial-approximation continuation}
\cfoot{\thepage}
\renewcommand{\headrulewidth}{0.4pt}

\titleformat{\section}{\Large\bfseries}{\thesection}{0.75em}{}
\titleformat{\subsection}{\large\bfseries}{\thesubsection}{0.75em}{}
\titleformat{\subsubsection}{\normalsize\bfseries}{\thesubsubsection}{0.75em}{}

\setlist[itemize]{leftmargin=2em,itemsep=0.3em,topsep=0.45em}
\setlist[enumerate]{leftmargin=2.2em,itemsep=0.3em,topsep=0.45em}
\setlist[description]{leftmargin=3.6cm,labelsep=0.5cm,itemsep=0.35em,topsep=0.45em}
\setlength{\parindent}{0pt}
\setlength{\parskip}{0.55em}
\setlength{\emergencystretch}{4em}

\newtheorem{theorem}{Theorem}[section]
\newtheorem{proposition}[theorem]{Proposition}
\newtheorem{lemma}[theorem]{Lemma}
\newtheorem{corollary}[theorem]{Corollary}
\newtheorem{conjecture}[theorem]{Conjecture}
\newtheorem{criterion}[theorem]{Criterion}
\newtheorem{question}[theorem]{Question}
\theoremstyle{definition}
\newtheorem{definition}[theorem]{Definition}
\newtheorem{observation}[theorem]{Observation}
\newtheorem{experiment}[theorem]{Experiment}
\newtheorem{warning}[theorem]{Caution}
\theoremstyle{remark}
\newtheorem{remark}[theorem]{Remark}

\newcommand{\N}{\mathbb N}
\newcommand{\Nzero}{\mathbb N_0}
\newcommand{\Z}{\mathbb Z}
\newcommand{\Q}{\mathbb Q}
\newcommand{\R}{\mathbb R}
\newcommand{\C}{\mathbb C}
\newcommand{\Bsem}{\mathcal B}
\newcommand{\Lsem}{\Lambda}
\newcommand{\cP}{\mathcal P}
\newcommand{\cL}{\mathcal L}
\newcommand{\cE}{\mathcal E}
\newcommand{\Int}{\operatorname{Int}}
\newcommand{\capc}{\operatorname{cap}}
\newcommand{\dist}{\operatorname{dist}}
\newcommand{\sgn}{\operatorname{sgn}}
\newcommand{\den}{\operatorname{den}}
\newcommand{\arsinh}{\operatorname{arsinh}}
\newcommand{\supp}{\operatorname{supp}}
\newcommand{\AEname}{Alaoglu--Erd\H{o}s}
\newcommand{\thetaq}{\vartheta}
\newcommand{\fall}[2]{\left(#1\right)_{\!\underline{#2}}}
\newcommand{\qbinom}[3]{\genfrac{[}{]}{0pt}{}{#1}{#2}_{#3}}
\newcommand{\statusnew}{\textcolor{proofgreen}{\textbf{[proved in this continuation]}}}
\newcommand{\statusclassical}{\textcolor{classicalgray}{\textbf{[classical]}}}
\newcommand{\statusinherited}{\textcolor{openblue}{\textbf{[inherited from the unified report]}}}
\newcommand{\statusexperiment}{\textcolor{computeviolet}{\textbf{[exact or numerical experiment]}}}
\newcommand{\statusopen}{\textcolor{openblue}{\textbf{[open target]}}}

\newenvironment{statusbox}[1]
{\begin{center}\begin{minipage}{0.94\textwidth}\raggedright\hrule\vspace{0.45em}\textbf{Status:} #1\par\vspace{0.35em}}
{\vspace{0.35em}\hrule\end{minipage}\end{center}}

\newenvironment{resultbox}
{\begin{center}\begin{minipage}{0.94\textwidth}\hrule\vspace{0.45em}}
{\vspace{0.35em}\hrule\end{minipage}\end{center}}

\lstdefinestyle{pythonstyle}{
  language=Python,
  basicstyle=\ttfamily\scriptsize,
  backgroundcolor=\color{codegray},
  frame=single,
  rulecolor=\color{rulegray},
  numbers=left,
  numberstyle=\tiny\color{rulegray},
  numbersep=8pt,
  showstringspaces=false,
  breaklines=true,
  columns=fullflexible,
  keepspaces=true,
  tabsize=4
}

\begin{document}
\pagenumbering{gobble}

\begin{titlepage}
\centering
\vspace*{0.55cm}
{\Huge\bfseries The \AEname\par}
\vspace{0.18cm}
{\Huge\bfseries Integer-Exponent Conjecture\par}
\vspace{0.65cm}
{\LARGE Polynomial Approximation, Integer-Valued Interpolation,\par}
\vspace{0.12cm}
{\LARGE and the Arithmetic Unit Barrier\par}
\vspace{0.85cm}
{\large A nonduplicative continuation of the supplied unified report\par}
\vfill
\begin{minipage}{0.91\textwidth}
\small
\textbf{Executive verdict.}
No unconditional proof of the conjecture is obtained, and no published resolution was found
in the literature checked through August 21, 2026.  The continuation does, however, extract a
new approximation-theoretic core with several exact results.  The sharp Bernstein-ellipse
optimization shows that a classical blockwise polynomial approximation can have enough nodes
only for the smallest hypothetical output $M=2^x=3$; the optimal block ratio is
$3+2\sqrt2$, and that entire degree-versus-node geometry cannot reach any $M\ge5$ before coefficient
arithmetic is considered.  A dyadic Newton-interpolation certificate then gives short exact
proofs that $M^\vartheta$ is nonintegral for $M=3,5,6,7$, where
$\vartheta=\log_2 3$.  A discrete minimax theorem identifies the real best error as a
cofactor quotient and shows exactly how clearing its denominator raises the residual to a
nonzero integer.  Finally, a $q$-Newton/Gaussian-binomial construction gives canonical
integer-valued interpolants on geometric progressions and an exact extrapolation formula;
its first omitted value grows rather than shrinks.  These results isolate the remaining task:
construct Bernstein-quality approximants in the sparse ring
$\operatorname{Int}(\langle2,3\rangle,\mathbb Z)$, or find an equivalent source of
nonarchimedean cancellation.
\end{minipage}
\vfill
{\large August 21, 2026\par}
\end{titlepage}

\pagenumbering{roman}

\begin{abstract}
Write $\thetaq=\log 3/\log 2$.  The \AEname\ conjecture asserts that
$2^x,3^x\in\Z$ for real $x$ implies $x\in\Z$; equivalently, an integer $M$ with
$M^{\thetaq}\in\Z$ must be a power of two.  The supplied unified report already develops
machine-checked structural reductions, exact conditional semigroup descriptions, determinant
methods, sparse interpolation, finite verification, and the reduction to the four
exponentials frontier.  This continuation does not repeat that material.  It asks instead
whether approximation theory can convert the many conditional integral values into a
subunit error and hence an exact contradiction.

The first result is an exact archimedean optimization.  If a hypothetical counterexample has
$M=2^\beta$, then the $3$-smooth nodes in $[X,cX]$ have asymptotic density
$\frac{\log c}{\log2\log3}\log X$, while best degree-$d$ polynomial approximation to
$t^\beta$ on that block requires
$d\sim\beta\log X/\log\rho(c)$, where
$\rho(c)=(\sqrt c+1)/(\sqrt c-1)$.  The efficiency
$\log c\log\rho(c)$ has the exact maximum
$2\log^2(1+\sqrt2)$, attained at $c=3+2\sqrt2$.  Consequently this scheme can beat the
node threshold only when
$M<\exp(2\log^2(1+\sqrt2)/\log3)=4.113124586\ldots$; among non-powers of two, only $M=3$
remains.  The second result explains why even that case is not settled by the real minimax
polynomial: normalization to a compact interval introduces a denominator cost of order
$X^d$, overwhelming the geometric gain.  Markov resolution gives the same diagnosis in a
different language: $O(\log X)$ arithmetic nodes are too coarse by a factor of $\log X$ to
norm degree $O(\log X)$ polynomials.

On the arithmetic side, this continuation proves a dyadic interpolation certificate.  If the
rational polynomial interpolating $T^{\thetaq}$ at $d+1$ powers of two comes within the
reciprocal of its reduced denominator at an integer $M$, then $M^{\thetaq}$ is not an
integer.  Quadratic instances give exact exclusions for $M=3,5,6,7$.  A general discrete
Chebyshev duality theorem expresses the best real uniform error on $d+2$ nodes as
$|\det V|$ divided by the sum of the adjacent cofactors.  At arithmetic nodes, clearing the
minimax denominator produces an integer-coefficient alternant whose error is the nonzero
integer $\det V$; this is the exact ``unit barrier.''  A separate $q$-Newton theorem constructs
integer-valued interpolation polynomials on $\{q^n\}$, with Gaussian-binomial basis and first
extrapolation error $\prod_{i=0}^d(r-q^i)$.  The formula makes clear why one-axis arithmetic
interpolation is plentiful but cross-axis control is absent.

The report closes by formulating precise sufficient conditions in the ring of
integer-valued polynomials on the $3$-smooth semigroup, by connecting them with Bhargava
$p$-orderings and generalized factorials, and by auditing rational, Pad\'e,
Hermite--Pad\'e, orthogonal-polynomial, and probabilistic-rounding approaches.  Recent work
on matrix coefficients, exponential-product transcendence, discrete Chebyshev systems, and
integer-valued polynomial bases sharpens the available language but does not remove the
rank-two four-exponentials obstruction.
\end{abstract}

\tableofcontents
\clearpage
\pagenumbering{arabic}

\section{Executive answer, scope, and claim ledger}
\label{sec:executive}

\subsection{What this continuation does not repeat}

The supplied unified report is already a book-length synthesis.  It contains the reduction to
four exponentials, the exact conditional free-monoid structure of all solutions, a
power-primitive odd generator, block counting and equidistribution, arbitrary-node and
semigroup determinants, $p$-adic valuation budgets, HCIZ bounds, sparse power-curve
interpolation, matrix-coefficient reformulations, boundary-dilation defects, finite scans,
and a detailed Lean inventory.  Repeating those arguments would obscure rather than extend
them.

This document uses only the following inherited facts.

\begin{enumerate}
\item A hypothetical noninteger solution is irrational and in fact transcendental.
\item If $\beta$ is such a solution and
      \[
        M=2^\beta,\qquad N=3^\beta,
      \]
      then $M,N$ are positive integers and
      \[
        (2^a3^b)^\beta=M^aN^b\in\Z
        \qquad(a,b\in\Nzero).
      \]
\item The functions $t^{\lambda_0},\ldots,t^{\lambda_s}$ with distinct real exponents form
      an extended complete Chebyshev system on $(0,\infty)$; a nonzero linear combination has
      at most $s$ positive zeros, counted without multiplicity for the applications here.
\item The conjecture is a special $2\times2$ case of the four exponentials problem, and the
      known six exponentials theorem does not specialize far enough.
\end{enumerate}

Everything else is rederived or newly developed below.

\subsection{Claim ledger}

\begin{longtable}{p{0.30\textwidth}p{0.19\textwidth}p{0.43\textwidth}}
\toprule
Statement & Status & Role \\
\midrule
\endfirsthead
\toprule
Statement & Status & Role \\
\midrule
\endhead
Bernstein root asymptotic for $t^\beta$ & \statusclassical & Determines the best possible
archimedean exponential rate on a fixed multiplicative block. \\
Exact maximization of $\log c\log\rho(c)$ & \statusnew & Shows that classical block
polynomial approximation has enough nodes only for $M=3$. \\
Arithmetic subunit locking on $3$-smooth nodes & \statusnew & Converts an
integer-valued approximant with error $<1$ into too many exact zeros. \\
Integer-part shift has no asymptotic gain & \statusnew & Shows that replacing $\beta$ by
its fractional part merely changes the arithmetic threshold from $1$ to $X^{-\lfloor\beta\rfloor}$. \\
Scaling/denominator tax & \statusnew & Proves that naive compact-interval
normalization costs $X^d$ when coefficients are cleared. \\
Dyadic interpolation certificate & \statusnew & Gives exact finite certificates from
rational interpolation at powers of two. \\
Exclusions $M=3,5,6,7$ & \statusnew & Independent symbolic consequences of the
certificate theorem; much weaker in range than inherited finite scans. \\
Discrete minimax cofactor identity & \statusnew & Gives the exact real best error and the
alternation vector on arbitrary arithmetic nodes. \\
Arithmetic alternant after denominator clearing & \statusnew & Identifies the residual
as the nonzero integer determinant, explaining the unit barrier. \\
$q$-Newton/Gaussian-binomial interpolation & \statusnew & Constructs integer-valued
interpolants on geometric progressions and gives exact extrapolation errors. \\
Large zero-free computational ranges & \statusinherited & Not rerun or enlarged here. \\
A complete proof of the conjecture & \statusopen & No such proof is claimed. \\
\bottomrule
\end{longtable}

\subsection{The new bottom line}

Approximation theory supplies an unexpectedly sharp archimedean opportunity.  For the
smallest hypothetical generator $M=3$, best polynomial approximation has enough degrees of
freedom to become subunit on more arithmetic nodes than its zero count permits.  This is not
merely a rough comparison: the optimal block and the exact threshold can be calculated.

The obstruction is arithmetic.  The best real approximant has coefficients with no reason to
take integral values on the conditional node set.  Clearing denominators changes a small
normalized error into a large integral residual.  The central problem may therefore be stated
without transcendence terminology:

\begin{resultbox}
Construct, for a nonintegral $\beta$ with $2^\beta,3^\beta\in\Z$, a degree-$d$ polynomial
that is integer-valued on $\langle2,3\rangle$, has error $<1$ from $t^\beta$ on at least
$d+2$ positive nodes, and retains Bernstein-quality size.
\end{resultbox}

The report proves that the real and arithmetic sides separately have the right scale in the
edge case $M=3$, but no construction currently satisfies both.

\section{Minimal setup and the approximation dictionary}
\label{sec:setup}

\subsection{Two dual normalizations}

Set
\[
  \thetaq:=\frac{\log3}{\log2}=1.584962500721156\ldots.
\]
There are two equivalent ways to organize a counterexample.

\paragraph{Exponent normalization.}
Assume that $\beta\notin\Z$ and
\[
  M=2^\beta\in\N,\qquad N=3^\beta\in\N.
\]
Then
\[
  \beta=\log_2M,\qquad M^{\thetaq}=N.
\]
The fixed arithmetic node set
\[
  \Bsem:=\{2^a3^b:a,b\in\Nzero\}
\]
satisfies $q^\beta\in\Z$ for every $q\in\Bsem$.

\paragraph{Power-curve normalization.}
Fix $\thetaq$ and regard $M$ as a positive integer.  The conjecture is equivalent to
\[
  M^{\thetaq}\in\Z\quad\Longrightarrow\quad M=2^m\text{ for some }m\in\Nzero.
\]
The first normalization is better for approximation on the fixed node set $\Bsem$; the
second is better for interpolation at the known points
$(2^m,3^m)$ of the curve $y=x^{\thetaq}$.

\subsection{The locking principle}

The basic mechanism is an integer-versus-small-real comparison.

\begin{lemma}[Chebyshev zero bound]
\label{lem:power-cheb}
Let $\lambda_0<\cdots<\lambda_s$ be real.  If
\[
  F(t)=\sum_{j=0}^s c_jt^{\lambda_j}
\]
is not identically zero, then $F$ has at most $s$ distinct zeros on $(0,\infty)$.
\end{lemma}

\begin{proof}
Put $t=e^u$.  Then $F(e^u)=\sum_jc_je^{\lambda_ju}$ is an exponential polynomial with
strictly ordered real frequencies.  Its Wronskian equals a nonzero Vandermonde factor times a
positive exponential.  Hence the functions form an extended complete Chebyshev system.
\end{proof}

\begin{theorem}[Arithmetic subunit locking]
\label{thm:locking}
Assume $q^\beta\in\Z$ for every $q$ in a finite set $E\subset(0,\infty)\cap\Z$.  Let
$P\in\R[T]$ have degree at most $d$ and satisfy $P(q)\in\Z$ for every $q\in E$
(for example, $P\in\operatorname{Int}(E,\Z)$).  If
\[
  \max_{q\in E}|q^\beta-P(q)|<1,
\]
then $q^\beta=P(q)$ for every $q\in E$.  Consequently, if $\beta\notin\Nzero$ and
$\#E\ge d+2$, no such $P$ exists.
\end{theorem}

\begin{proof}
For each $q\in E$, both $q^\beta$ and $P(q)$ are integers.  A difference of absolute value
less than one must vanish.  If there are $d+2$ zeros, the nonzero generalized polynomial
$q^\beta-P(q)$ contradicts Lemma~\ref{lem:power-cheb} applied to
$1,t,\ldots,t^d,t^\beta$.
\end{proof}

The same argument does not require ordinary powers.

\begin{theorem}[Quantized M\"untz approximation]
\label{thm:muntz-lock}
Assume $M=2^\beta,N=3^\beta\in\Z$ and define
\[
  \Lsem_\beta:=\{n+k\beta:n,k\in\Nzero\}.
\]
For every $q\in\Bsem$ and $\lambda\in\Lsem_\beta$, one has $q^\lambda\in\Z$.
Let $\lambda_1,\ldots,\lambda_s,\mu$ be distinct elements of $\Lsem_\beta$.  For any finite
$E\subset\Bsem$ with $\#E\ge s+1$,
\[
 \inf_{c_1,\ldots,c_s\in\Z}
 \max_{q\in E}\left|q^\mu-\sum_{j=1}^sc_jq^{\lambda_j}\right|\ge1.
\]
Likewise, every nonzero integer-coefficient M\"untz polynomial with $s$ terms has supremum at
least one on any $s$ distinct nodes of $\Bsem$.
\end{theorem}

\begin{proof}
Integrality follows from
\[
 (2^a3^b)^{n+k\beta}=(2^a3^b)^n(M^aN^b)^k.
\]
If the displayed maximum were below one, all residuals would vanish.  The residual has at
most $s$ positive zeros by Lemma~\ref{lem:power-cheb}, whereas $E$ has at least $s+1$ nodes.
The homogeneous statement is identical with one fewer term.
\end{proof}

\begin{remark}
Theorem~\ref{thm:muntz-lock} is the approximation-theoretic dual of an integer generalized
Vandermonde determinant.  The determinant language is already developed in the supplied
report; the new point here is that the lower bound is exactly the arithmetic unit $1$, while
the corresponding real minimax error may be exponentially smaller.
\end{remark}

\subsection{Subtracting the integer part does not help}

One might hope to reduce a large $\beta$ to its fractional part before approximating.  The
arithmetic threshold changes by exactly the amount that cancels this advantage.

\begin{lemma}[Integer-shift invariance of the degree scale]
\label{lem:integer-shift}
Write $\beta=m+\delta$ with $m\in\Nzero$ and $0<\delta<1$.  For $q\in\Bsem$ and
$P\in\Z[T]$,
\[
 q^m\bigl(q^\delta-P(q)\bigr)=q^\beta-q^mP(q)\in\Z.
\]
Thus a sufficient pointwise locking condition is
\[
 |q^\delta-P(q)|<q^{-m}.
\]
On a block $[X,cX]$, best degree-$d$ real polynomial approximation to $t^\delta$ has
archimedean scale $X^\delta\rho(c)^{-d}$, so reaching the threshold $X^{-m}$ still requires
\[
 d\ge \frac{\beta\log X}{\log\rho(c)}+o(\log X).
\]
\end{lemma}

\begin{proof}
The integrality identity is immediate.  The approximation statement follows from the
Bernstein root asymptotic proved in Section~\ref{sec:bernstein}: the condition
$X^\delta\rho(c)^{-d}<X^{-m+o(1)}$ is equivalent to the displayed degree bound.
\end{proof}

This simple cancellation is a recurring warning: lowering the analytic exponent usually
introduces an exactly compensating denominator.

\section{Recent literature and the relevant frontier}
\label{sec:literature}

\subsection{Four exponentials and matrix coefficients}

Waldschmidt's 2023 survey places the question ``$2^t,3^t\in\Z$ implies $t\in\Nzero$?''
explicitly inside the four exponentials problem and records that both remain open
\cite{Waldschmidt2023}.  Dasgupta and Kakde subsequently formulated the Matrix Coefficient
Conjecture for matrices of logarithms of algebraic numbers and proved a substantial converse
direction linking it with four exponentials and structural rank
\cite{DasguptaKakde2024,DasguptaKakde2026}.  The supplied unified report already specializes
that framework to the present $2\times2$ matrix.  Approximation theory does not bypass this
frontier automatically: its coefficient lattice is another form of the same rank-collapse
problem.

Massold's 2025 preprint gives new lower bounds for transcendence degrees of fields generated
by exponentials of products and reduces a stronger estimate to a torus-intersection
conjecture \cite{Massold2025}.  The parameter range does not supply the missing $2\times2$
case here.  The importance of the work is diagnostic: even contemporary advances naturally
stop at a geometric intersection statement when the product matrix is critical.

\subsection{Hermite--Pad\'e and explicit transcendence measures}

Fischler and Rivoal's 2025 preprint develops a new explicit transcendence measure for values
of the exponential function.  The construction optimizes an accessory parameter in a Siegel
determinant built from Hermite--Pad\'e approximants to powers of $\exp$
\cite{FischlerRivoal2025}.  This is very close in spirit to what one would like here: an
auxiliary function with unusually small archimedean value and controlled arithmetic
coefficients.

The hypotheses, however, are materially different.  Classical $E$-function and
Hermite--Pad\'e machinery obtains arithmetic information from differential equations over
algebraic function fields and evaluation at algebraic points.  In the present problem the
natural functions are
\[
 z\longmapsto e^{z\log2},\qquad z\longmapsto e^{z\log3},
\]
whose differential coefficients $\log2,\log3$ are transcendental, and the hypothetical
evaluation point $z=\beta$ is transcendental.  Faverjon and Adamczewski's recent algebraic
independence measures for $E$- and $M$-functions likewise emphasize algebraic systems and
algebraic evaluation points \cite{AdamczewskiFaverjon2025}.  Section~\ref{sec:hermite-pade}
turns this mismatch into an explicit failure audit.

\subsection{Discrete Chebyshev systems and integer-valued polynomials}

Gorbachev, Ivanov, and Tikhonov developed a discrete Chebyshev alternance theorem and a
Remez-type algorithm for restrictions of Chebyshev systems to ordered discrete sets
\cite{GorbachevIvanovTikhonov2025}.  Their framework is directly relevant to the cofactor
minimax identity in Section~\ref{sec:minimax}: the conditional arithmetic nodes form a
strictly ordered discrete set and the power functions form a Chebyshev system.

On the arithmetic side, Berger and Sprang connected regular bases of integer-valued
polynomial rings with $p$-adic Fourier theory \cite{BergerSprang2025}.  Lagarias and Yangjit's 2025 work on $B$-orderings extends Bhargava's ordering machinery for generalized
factorials \cite{LagariasYangjit2025,Bhargava1997}.  These developments do not solve the
\AEname\ problem, but they provide exactly the language needed for the open construction in
$\Int(\Bsem,\Z)$.

\subsection{Integer Chebyshev and potential theory}

The integer Chebyshev problem studies small polynomials with integer coefficients on compact
sets, with logarithmic capacity governing the first approximation scale.  The monic variant
was developed in detail by Borwein, Pinner, and Pritsker
\cite{BorweinPinnerPritsker2003}; Pritsker's work surveys the interaction of capacity,
algebraic integers, and integer coefficients \cite{Pritsker2005}.  Our problem is
inhomogeneous and the relevant ring is integer-valued on a sparse semigroup rather than
$\Z[T]$.  Nevertheless, potential theory explains why normalization to $[1,c]$ is
archimedeanly optimal and arithmetically dangerous.

\begin{statusbox}
The literature review above is \statusclassical.  No source located through August 21, 2026
claims a proof of the \AEname\ conjecture or of the general four exponentials conjecture.
The new theorems below are independent paper proofs and are not claimed to be machine checked.
\end{statusbox}

\section{Bernstein ellipses and the exact block-efficiency theorem}
\label{sec:bernstein}

\subsection{Best approximation on a multiplicative block}

For $c>1$, $\alpha>0$ with $\alpha\notin\Nzero$, and $d\in\Nzero$, define
\[
 E_d(\alpha;c):=
 \inf_{P\in\R[T],\,\deg P\le d}\ 
 \max_{1\le t\le c}|t^\alpha-P(t)|.
\]
The affine map from $[1,c]$ to $[-1,1]$ sends the branch point $t=0$ to
\[
 s_0=-\frac{c+1}{c-1}<-1.
\]
The Bernstein ellipse through $s_0$ has parameter
\begin{equation}
 \rho(c)=|s_0|+\sqrt{s_0^2-1}
 =\frac{\sqrt c+1}{\sqrt c-1}.
 \label{eq:rho}
\end{equation}

\begin{theorem}[Bernstein root asymptotic]
\label{thm:bernstein-root}
For fixed $c>1$ and positive nonintegral $\alpha$,
\[
 \lim_{d\to\infty} E_d(\alpha;c)^{1/d}=\rho(c)^{-1}.
\]
Consequently, on $[X,cX]$,
\[
 \inf_{\deg P\le d}\max_{X\le t\le cX}|t^\alpha-P(t)|
 =X^\alpha E_d(\alpha;c),
\]
and subunit real error requires
\begin{equation}
 d\ge \frac{\alpha\log X}{\log\rho(c)}+o(\log X).
 \label{eq:degree-need}
\end{equation}
\end{theorem}

\begin{proof}
The function $t^\alpha$, with the principal branch, is holomorphic in the slit plane
$\C\setminus(-\infty,0]$.  Under the affine map, it is holomorphic inside every Bernstein
ellipse of parameter smaller than $\rho(c)$, and the branch point at $0$ prevents analytic
continuation through the ellipse of parameter $\rho(c)$.  Bernstein's converse theorem for
geometric polynomial approximation gives the root asymptotic.  The scaling identity follows
by writing $t=Xu$ and absorbing $X^\alpha$ into the coefficients.
\end{proof}

\subsection{How many arithmetic nodes are available?}

Let
\[
 B(X,c):=\#\bigl(\Bsem\cap[X,cX]\bigr).
\]
Put $a=\log2$, $b=\log3$, $L=\log X$, and $h=\log c$.  The condition
$2^m3^n\in[X,cX]$ is
\[
 L\le am+bn\le L+h.
\]
Since $a/b$ is irrational, the boundary phases are equidistributed.

\begin{proposition}[Thin-strip count]
\label{prop:thin-strip}
For every fixed $c>1$,
\begin{equation}
 B(X,c)=\frac{\log c}{\log2\log3}\log X+o(\log X)
 \qquad(X\to\infty).
 \label{eq:node-count}
\end{equation}
\end{proposition}

\begin{proof}[Proof sketch]
Sum over $n\le(L+h)/b$.  For each $n$, the admissible $m$ lie in an interval of fixed length
$h/a$ whose endpoints have fractional parts obtained from the irrational rotation
$n(b/a)$.  Weyl equidistribution gives average occupancy $h/a$; the number of $n$ is
$L/b+O(1)$.  The resulting main term is $hL/(ab)$ and the discrepancy is $o(L)$.
\end{proof}

The supplied report proves stronger blockwise equidistribution statements conditionally for
candidate semigroups.  Proposition~\ref{prop:thin-strip} is the fixed-node analogue needed
here.

\subsection{The necessary degree-versus-node inequality}

Assume a counterexample $\beta$ and suppose, optimistically, that one can find a degree-$d$
polynomial integer-valued on $\Bsem$ whose error from $t^\beta$ is below one throughout
$[X,cX]$.  Theorem~\ref{thm:locking} needs $B(X,c)\ge d+2$.  Comparing
\eqref{eq:degree-need} and \eqref{eq:node-count} gives the leading-order necessary inequality
\begin{equation}
 \frac{\log c}{\log2\log3}\ge
 \frac{\beta}{\log\rho(c)}.
 \label{eq:raw-efficiency}
\end{equation}
Because $M=2^\beta$, this is
\begin{equation}
 \boxed{\quad
   \log c\,\log\rho(c)\ge\log3\,\log M.
 \quad}
 \label{eq:efficiency-condition}
\end{equation}
A strict inequality is needed for a positive asymptotic degree margin; equality would require
lower-order analysis.

The left side depends only on block geometry.  It can be optimized exactly.

\subsection{Exact optimization}

\begin{theorem}[Optimal multiplicative block]
\label{thm:optimal-block}
For $c>1$, set
\[
 \Psi(c):=\log c\,\log\frac{\sqrt c+1}{\sqrt c-1}.
\]
Then
\begin{equation}
 \max_{c>1}\Psi(c)=2\log^2(1+\sqrt2),
 \label{eq:psi-max}
\end{equation}
and the unique maximizer is
\begin{equation}
 c_*=3+2\sqrt2=(1+\sqrt2)^2.
 \label{eq:cstar}
\end{equation}
At this point $\rho(c_*)=1+\sqrt2$.
\end{theorem}

\begin{proof}
Write
\[
 y=\tfrac12\log c,
 \qquad
 u=\log\rho(c)=\log\coth(y/2).
\]
The identity $\tanh(y/2)=e^{-u}$ gives
\[
 \sinh y\,\sinh u=1.
\]
We claim that $yu\le\arsinh(1)^2$.  The function
\[
 z\longmapsto \log\sinh(e^z)
\]
is strictly convex, because with $x=e^z$ its second derivative is
\[
 x\coth x-x^2\operatorname{csch}^2x
 =\frac{x(\sinh x\cosh x-x)}{\sinh^2x}>0.
\]
Applying midpoint convexity at $\log y$ and $\log u$ yields
\[
 \sinh\sqrt{yu}
 \le\sqrt{\sinh y\sinh u}=1.
\]
Hence $\sqrt{yu}\le\arsinh(1)=\log(1+\sqrt2)$.  Equality forces $y=u$,
so $\sinh^2y=1$, whence $y=u=\arsinh(1)$.  Finally
$\Psi(c)=2yu$, giving \eqref{eq:psi-max}, and
$c=e^{2y}=(1+\sqrt2)^2$.
\end{proof}

\begin{corollary}[The generator-three frontier]
\label{cor:M3-frontier}
A blockwise Bernstein-rate polynomial strategy can satisfy the leading-order
degree-versus-node inequality only if
\begin{equation}
 M\le\exp\!\left(\frac{2\log^2(1+\sqrt2)}{\log3}\right)
 =4.113124586019410\ldots.
 \label{eq:Mcritical}
\end{equation}
A positive asymptotic margin requires the strict inequality.  Among positive integers that
are not powers of two, only $M=3$ qualifies.
\end{corollary}

\begin{proof}
Combine \eqref{eq:efficiency-condition} with Theorem~\ref{thm:optimal-block}.  The numerical
value follows by direct evaluation.  The only integer $M$ with $2<M<4.113\ldots$ is $3$;
$M=4$ corresponds to the integral exponent $\beta=2$.
\end{proof}

\begin{table}[ht]
\centering
\small
\begin{tabular}{rrrrr}
\toprule
$M$ & $\beta=\log_2M$ & node coefficient & degree coefficient & margin \\
\midrule
$3$ & $1.5849625$ & $2.3148354$ & $1.7982868$ & $+0.5165486$ \\
$4$ & $2.0000000$ & $2.3148354$ & $2.2691853$ & $+0.0456501$ \\
$5$ & $2.3219281$ & $2.3148354$ & $2.6344426$ & $-0.3196072$ \\
$6$ & $2.5849625$ & $2.3148354$ & $2.9328795$ & $-0.6180441$ \\
$7$ & $2.8073549$ & $2.3148354$ & $3.1852043$ & $-0.8703689$ \\
$9$ & $3.1699250$ & $2.3148354$ & $3.5965736$ & $-1.2817382$ \\
\bottomrule
\end{tabular}
\caption{Asymptotic coefficients of $\log X$ at the optimal block $c_*=3+2\sqrt2$.
The node coefficient is $\log c_*/(\log2\log3)$; the degree coefficient is
$\beta/\log(1+\sqrt2)$.  $M=4$ is shown only as a calibration.}
\label{tab:efficiency}
\end{table}

\subsection{Interpretation}

Corollary~\ref{cor:M3-frontier} is both encouraging and restrictive.

\begin{itemize}
\item It is encouraging because the archimedean side is not uniformly too weak.  For $M=3$
      it has a genuine positive asymptotic margin.
\item It is restrictive because no sharpening of constants inside ordinary polynomial
      approximation can make this particular block-count method reach $M\ge5$.  A different
      function space, a different arithmetic source of nodes, or a nonarchimedean gain is
      necessary.
\item It identifies the exact place where an integer-valued approximation theorem would have
      the highest leverage: the ring $\Int(\Bsem,\Z)$, exponent
      $\beta=\thetaq$, and block ratio $3+2\sqrt2$.
\end{itemize}

The next section explains why the real minimax polynomial cannot simply be rounded or scaled
into that ring.

\section{Why the real minimax polynomial does not become arithmetic}
\label{sec:integerization}

\subsection{The normalization and denominator tax}

A best approximant on $[1,c]$ has real coefficients.  To use
Theorem~\ref{thm:locking}, one must turn it into a polynomial taking integer values at the
arithmetic nodes.  The most obvious conversion is fatal.

\begin{proposition}[Naive coefficient clearing]
\label{prop:clearing-tax}
Assume $X\in\Bsem$ and put $A=X^\beta\in\Z$.  Let
\[
 p(u)=\sum_{j=0}^da_ju^j\in\Q[u],
 \qquad Da_j\in\Z
\]
for a positive integer $D$.  Define
\[
 \widetilde P_X(t):=DX^dA\,p(t/X)\in\Z[t].
\]
Then
\begin{equation}
 \max_{X\le t\le cX}
 \bigl|DX^dt^\beta-\widetilde P_X(t)\bigr|
 =DX^{d+\beta}\max_{1\le u\le c}|u^\beta-p(u)|.
 \label{eq:clearing-error}
\end{equation}
If $p$ has Bernstein-rate error $\rho(c)^{-d+o(d)}$, the logarithm of the right side is
\[
 d\log X-d\log\rho(c)+\beta\log X+\log D+o(d).
\]
For every choice $d\asymp\log X$, the positive term $d\log X\asymp(\log X)^2$ overwhelms the
Bernstein gain $d\log\rho(c)\asymp\log X$.
\end{proposition}

\begin{proof}
The coefficient of $t^j$ in $DX^dA p(t/X)$ is
$DAa_jX^{d-j}\in\Z$.  The identity \eqref{eq:clearing-error} follows after setting $t=Xu$.
The asymptotic comparison is immediate.
\end{proof}

\begin{warning}
Proposition~\ref{prop:clearing-tax} is a no-go theorem only for the naive operation
``approximate on $[1,c]$, rescale, clear all coefficient denominators.''  It does not rule out
polynomials that are integer-valued on $\Bsem$ without having integer coefficients.  That
larger ring is the principal open direction in Section~\ref{sec:int-valued}.
\end{warning}

The same tax appears if one normalizes several $3$-smooth nodes by a common node $X$.  Ratios
such as $2^a3^b/X$ have numerator and denominator of order $X$ when the exponent vectors move
in opposite directions.  A common coefficient denominator for a degree-$d$ polynomial in
those ratios can therefore be of order $X^d$.  The obstruction is not numerical roundoff; it
is the arithmetic cost of changing scale.

\subsection{Capacity on the original scale}

Potential theory gives a complementary formulation.

\begin{proposition}[Monic Chebyshev norm on a block]
\label{prop:monic-capacity}
Let $I=[X,cX]$ and $d\ge1$.  Among monic real polynomials of degree $d$,
\begin{equation}
 \min_{P\text{ monic}}\|P\|_I
 =2\left(\frac{(c-1)X}{4}\right)^d.
 \label{eq:monic-cheb}
\end{equation}
In particular,
\[
 \capc(I)=\frac{(c-1)X}{4}.
\]
\end{proposition}

\begin{proof}
Scale the monic Chebyshev polynomial on $[-1,1]$.  Since $T_d$ has leading coefficient
$2^{d-1}$ and uniform norm one, the monic extremal polynomial on an interval of length
$L=(c-1)X$ has norm $L^d/2^{2d-1}=2(L/4)^d$.
\end{proof}

On the normalized interval $[1,c]$, capacity is fixed and geometric approximation is
possible.  On the original interval, capacity grows linearly with $X$.  Integer coefficients
live naturally on the original scale, whereas complex approximation lives naturally on the
normalized scale.  Proposition~\ref{prop:clearing-tax} is the coefficient-level expression
of this conflict.

The classical integer Chebyshev problem asks for a nonzero $P\in\Z[T]$ with small norm on a
fixed compact set.  Here one asks for an inhomogeneous approximation to $t^\beta$, and the
right arithmetic class is $\Int(\Bsem,\Z)$.  Capacity remains an indispensable lower-order
constraint, but it cannot alone see the sparse congruences of $\Bsem$.

\subsection{The Markov resolution deficit}

Could one control a polynomial on the full block from its values on the arithmetic nodes and
then use a real minimax construction?  Standard norming arguments fail by one logarithm.

For a finite set $E\subset I=[X,cX]$, define its fill distance
\[
 h(E,I):=\max_{t\in I}\min_{q\in E}|t-q|.
\]
Markov's inequality on an interval of length $L=(c-1)X$ gives
\[
 \|P'\|_I\le\frac{2d^2}{L}\|P\|_I
 \qquad(\deg P\le d).
\]
Therefore
\begin{equation}
 \|P\|_I
 \le \|P\|_E+\frac{2d^2h(E,I)}{L}\|P\|_I.
 \label{eq:markov-norming}
\end{equation}
When $2d^2h/L<1$, the nodes are a norming mesh with constant
$(1-2d^2h/L)^{-1}$.

\begin{proposition}[Arithmetic nodes are below Markov resolution]
\label{prop:markov-deficit}
Let $E_X=\Bsem\cap[X,cX]$, with fixed $c>1$.  If $d\asymp\log X$, then
\[
 \frac{2d^2h(E_X,[X,cX])}{(c-1)X}\longrightarrow\infty
\]
after passage to any sequence on which the implicit degree constant is positive.  More
precisely,
\[
 h(E_X,[X,cX])\ge
 \frac{(c-1)X}{2(B(X,c)+1)},
\]
so the Markov factor is at least $d^2/(B(X,c)+1)\asymp\log X$.
\end{proposition}

\begin{proof}
The $B(X,c)$ nodes split the interval into at most $B(X,c)+1$ gaps.  One gap has length at
least $L/(B(X,c)+1)$, and its midpoint is at distance at least half that length from every
node.  Proposition~\ref{prop:thin-strip} gives $B(X,c)\asymp\log X$.
\end{proof}

A degree-$d$ polynomial has boundary-layer resolution $L/d^2$.  The $3$-smooth nodes have
spacing scale $L/d$ when $d\asymp\log X$.  Thus a degree-$d$ polynomial can form a narrow peak
inside a typical arithmetic gap.  This does not prove that every discrete approximation
scheme fails, but it proves that the most common Markov-mesh transfer cannot supply the
missing estimate.

A measure-theoretic Remez inequality is even less adapted: a finite node set has measure zero.
Thickening each node to the Markov scale $L/d^2$ gives total measure only
$O(B(X,c)L/d^2)=O(L/\log X)$, and the resulting Remez constant grows roughly like a power of
$\log X$ raised to degree $d$.  Its logarithm is of order
$\log X\log\log X$, larger than the $O(\log X)$ Bernstein margin.

\subsection{Rounding coefficients is not a small perturbation}

Let $p(t)=\sum_{j=0}^da_jt^j$ approximate $t^\beta$ on $[X,cX]$.  Rounding each coefficient
to a nearby integer creates a pointwise perturbation bounded only by
\[
 \frac12\sum_{j=0}^d(cX)^j\asymp X^d.
\]
Using a Newton or Chebyshev basis reduces conditioning over the real numbers but does not by
itself make the basis integer-valued on $\Bsem$.  The correct rounding problem is a closest
vector problem in the evaluation lattice
\[
 \mathcal L_E=
 \bigl\{(P(q))_{q\in E}:P\in\Z[T],\ \deg P\le d\bigr\}\subset\Z^E,
\]
or, more flexibly, in the evaluation image of $\Int(\Bsem,\Z)$.  The cofactor identity of
Section~\ref{sec:minimax} computes the continuous optimum, but its distance to this lattice is
the unresolved quantity.

\section{Dyadic interpolation as an exact proof-producing certificate}
\label{sec:dyadic-cert}

The power-curve normalization supplies infinitely many exact rational interpolation data:
\[
 (2^n,(2^n)^{\thetaq})=(2^n,3^n),\qquad n\in\Nzero.
\]
A rational interpolant through such points can be compared with an integer target at another
integer $M$.  The reciprocal of its reduced denominator is an exact exclusion threshold.

\subsection{The certificate theorem}

\begin{theorem}[Dyadic interpolation certificate]
\label{thm:dyadic-certificate}
Let $0\le n_0<\cdots<n_d$ and let $I_{\mathbf n}\in\Q[T]$ be the unique polynomial of degree
at most $d$ satisfying
\[
 I_{\mathbf n}(2^{n_j})=3^{n_j}\qquad(0\le j\le d).
\]
Let $M\in\N$ be distinct from all $2^{n_j}$, and write
\[
 I_{\mathbf n}(M)=\frac ab
\]
in lowest terms with $b>0$.  If
\begin{equation}
 \left|M^{\thetaq}-\frac ab\right|<\frac1b,
 \label{eq:dyadic-cert}
\end{equation}
then $M^{\thetaq}\notin\Z$.
\end{theorem}

\begin{proof}
Suppose $M^{\thetaq}\in\Z$.  The difference
$M^{\thetaq}-a/b$ is rational.  It cannot vanish: otherwise
$T^{\thetaq}-I_{\mathbf n}(T)$ would have the $d+2$ distinct positive zeros
$2^{n_0},\ldots,2^{n_d},M$, contradicting Lemma~\ref{lem:power-cheb}.  A nonzero difference
between an integer and a reduced fraction with denominator $b$ has absolute value at least
$1/b$, contradicting \eqref{eq:dyadic-cert}.
\end{proof}

\begin{remark}
This is a genuine polynomial-approximation proof certificate, not a floating-point test.  The
interpolant and its denominator are exact.  The only analytic input is a certified interval
for $M^{\thetaq}$ narrow enough to verify \eqref{eq:dyadic-cert}.  Such an interval can be
proved by integer power comparisons using rational bounds for $\thetaq$.
\end{remark}

\subsection{Two quadratic interpolants}

Interpolation at $1,2,4$ gives
\begin{equation}
 I_{0,1,2}(T)=\frac{T^2+3T-1}{3}.
 \label{eq:I012}
\end{equation}
Interpolation at $2,4,8$ gives the scaled polynomial
\begin{equation}
 I_{1,2,3}(T)=\frac{T^2+6T-4}{4}.
 \label{eq:I123}
\end{equation}
Both identities are checked by substitution.

The following rational bounds are elementary:
\begin{equation}
 \frac{19}{12}<\thetaq<\frac85,
 \qquad
 \thetaq<\frac{27}{17}.
 \label{eq:theta-rational-bounds}
\end{equation}
Indeed,
\[
 3^{12}=531441>524288=2^{19},
 \qquad
 3^5=243<256=2^8,
\]
and
\[
 3^{17}=129140163<134217728=2^{27}.
\]

\begin{corollary}[Four short symbolic exclusions]
\label{cor:small-M}
One has
\[
 3^{\thetaq},\ 5^{\thetaq},\ 6^{\thetaq},\ 7^{\thetaq}\notin\Z.
\]
Equivalently, no \AEname\ counterexample can have
$2^\beta\in\{3,5,6,7\}$.
\end{corollary}

\begin{proof}
For $M=3$, \eqref{eq:I012} gives $I(3)=17/3$.  The exact comparisons
\[
 3^{31}=617673396283947>582622237229761=17^{12},
 \qquad
 3^8=6561<7776=6^5
\]
imply
\[
 \frac{17}{3}<3^{19/12}<3^{\thetaq}<3^{8/5}<6.
\]
Hence $|3^{\thetaq}-17/3|<1/3$, and
Theorem~\ref{thm:dyadic-certificate} applies.

For $M=5$, the same interpolant gives $I(5)=13$.  The integer comparisons
\[
 5^{19}=19073486328125>8916100448256=12^{12},
 \qquad
 5^8=390625<537824=14^5
\]
give $12<5^{\thetaq}<14$, so $|5^{\thetaq}-13|<1$.

For $M=6$, \eqref{eq:I123} gives $I(6)=17$.  Since
\[
 6^{19}=609359740010496>582622237229761=17^{12},
 \qquad
 6^8=1679616<1889568=18^5,
\]
we obtain $17<6^{\thetaq}<18$.

For $M=7$, \eqref{eq:I123} gives $I(7)=87/4$.  The lower comparison
\[
 7^{19}4^{12}=191461105130195933708288
 >188031682201497672618241=87^{12}
\]
and the upper comparisons in \eqref{eq:theta-rational-bounds}, together with
\[
 7^{27}=65712362363534280139543
 <8771994066994086373138432=22^{17},
\]
give
\[
 \frac{87}{4}<7^{19/12}<7^{\thetaq}<7^{27/17}<22.
\]
Thus $|7^{\thetaq}-87/4|<1/4$.
\end{proof}

\begin{table}[ht]
\centering
\small
\begin{tabular}{lllll}
\toprule
$M$ & interpolation nodes & $I(M)$ & certified interval for $M^{\thetaq}$ & threshold \\
\midrule
$3$ & $1,2,4$ & $17/3$ & $(17/3,6)$ & $1/3$ \\
$5$ & $1,2,4$ & $13$ & $(12,14)$ & $1$ \\
$6$ & $2,4,8$ & $17$ & $(17,18)$ & $1$ \\
$7$ & $2,4,8$ & $87/4$ & $(87/4,22)$ & $1/4$ \\
\bottomrule
\end{tabular}
\caption{Exact dyadic interpolation certificates.  The open interval has width at most twice
the threshold around the displayed interpolant value; equality with the interpolant is
excluded by the Chebyshev zero bound.}
\label{tab:small-certificates}
\end{table}

\subsection{What the certificate adds, and what it does not}

The inherited finite verification in the supplied report excludes a vastly larger range of
candidate generators.  Corollary~\ref{cor:small-M} is therefore not a record computation.
Its value is conceptual:

\begin{itemize}
\item it is a short symbolic proof based entirely on polynomial approximation and exact
      integer inequalities;
\item it makes the denominator threshold explicit;
\item it produces independently checkable certificates consisting only of a node list, a
      rational polynomial value, and a rational interval for $M^{\thetaq}$;
\item it generalizes immediately to any set of dyadic interpolation nodes.
\end{itemize}

For larger $M$, the interpolation error and the reduced denominator compete.  Moving the
nodes near $M$ improves real approximation but introduces powers of two into the denominator.
This is a finite version of Proposition~\ref{prop:clearing-tax}.

\subsection{A rational-interpolation variant}

Let $P,Q\in\Z[T]$ be coprime and suppose $P/Q$ interpolates $T^{\thetaq}$ at
$d+e+1$ dyadic nodes, with $\deg P\le d$ and $\deg Q\le e$.  If $Q(M)\ne0$ and
\begin{equation}
 |M^{\thetaq}Q(M)-P(M)|<1,
 \label{eq:rational-finite-cert}
\end{equation}
then $M^{\thetaq}\notin\Z$.  Indeed, an integral target would make the residual a nonzero
integer.  It is nonzero because
$T^{\thetaq}Q(T)-P(T)$ contains at most $d+e+2$ distinct powers and therefore at most
$d+e+1$ positive zeros, while equality at $M$ would add one zero beyond the interpolation
count.  Section~\ref{sec:rational} develops the asymptotic version and explains why the factor
$Q(M)$ is usually prohibitive.

\section{Discrete minimax duality and the exact unit barrier}
\label{sec:minimax}

Real best approximation on a finite node set has a closed determinant formula.  Its
arithmetic interpretation reveals the missing denominator saving.

\subsection{The cofactor alternant}

Let $q_0<\cdots<q_N$ be positive nodes and let
$\lambda_0<\cdots<\lambda_{N-1}<\mu$.  Put
\[
 A=(q_i^{\lambda_j})_{\substack{0\le i\le N\\0\le j<N}},
 \qquad
 f=(q_i^\mu)_{0\le i\le N},
 \qquad
 V=[A\ f].
\]
For $0\le i\le N$, let
\[
 C_i=(-1)^{i+N}\det A_{\widehat i},
\]
where $A_{\widehat i}$ is obtained by deleting row $i$.  Set
\[
 D=\det V,
 \qquad
 S=\sum_{i=0}^N|C_i|.
\]
Strict total positivity of the generalized Vandermonde matrix implies $D>0$ and that the
$C_i$ alternate in sign after a fixed global normalization.

\begin{theorem}[Exact discrete minimax error]
\label{thm:cofactor-minimax}
With the notation above,
\begin{equation}
 \min_{c\in\R^N}\|f-Ac\|_\infty=\frac{|D|}{S}.
 \label{eq:minimax-cofactor}
\end{equation}
There is a unique minimizer, and its residual vector is
\begin{equation}
 f-Ac_* = \frac{D}{S}\bigl(\sgn C_0,\ldots,\sgn C_N\bigr)^T.
 \label{eq:alternating-residual}
\end{equation}
\end{theorem}

\begin{proof}
The vector $C=(C_0,\ldots,C_N)^T$ spans the one-dimensional nullspace of $A^T$ by Laplace
expansion.  The $\ell^\infty$ distance from $f$ to $\operatorname{col}(A)$ has the dual
formula
\[
 \min_c\|f-Ac\|_\infty
 =\max_{A^Ty=0,\ \|y\|_1\le1}|y^Tf|.
\]
The feasible nullspace is one-dimensional, so the maximum is attained at
$y=C/S$.  Another Laplace expansion gives $C^Tf=D$.  Hence the optimum is $|D|/S$.

Let $s_i=\sgn C_i$.  The dual equality condition for $\ell^\infty$--$\ell^1$ duality forces
the primal residual to equal $\varepsilon s$ at every node, with
$\varepsilon=D/S$.  The augmented matrix $[A\ s]$ is nonsingular because its determinant is
$\sum_i|C_i|=S$, so this residual and the minimizer are unique.
\end{proof}

This is the finite-node form of Chebyshev alternation.  It also appears naturally in the
discrete theory of Chebyshev systems developed in \cite{GorbachevIvanovTikhonov2025}.

\subsection{Clearing the minimax denominator}

Suppose now that all entries of $A$ and $f$ are integers, as happens conditionally on
$\Bsem$ for exponents in $\Lsem_\beta$.  Then $D,C_i,S\in\Z$.

\begin{theorem}[Arithmetic alternant]
\label{thm:arithmetic-alternant}
Under the integral-entry hypothesis, let $s=(\sgn C_i)$.  There exists
$z\in\Z^N$ such that
\begin{equation}
 Sf-Az=D s.
 \label{eq:integer-alternant}
\end{equation}
Equivalently, the integer-coefficient generalized polynomial
\[
 G(t)=S t^\mu-\sum_{j=0}^{N-1}z_jt^{\lambda_j}
\]
takes the alternating values $D\,\sgn C_i$ at the nodes $q_i$.
In particular, the real minimax residual $D/S$ becomes the nonzero integer residual $D$
after the natural denominator is cleared.
\end{theorem}

\begin{proof}
Solve $Az=Sf-Ds$.  By \eqref{eq:alternating-residual}, the real solution is $z=Sc_*$.  Since
$\det[A\ s]=S$,
\[
 \binom{z}{D}=[A\ s]^{-1}Sf=\operatorname{adj}([A\ s])f.
\]
The adjugate and $f$ are integral, so every component of $z$ is an integer.
\end{proof}

\begin{remark}
Theorem~\ref{thm:arithmetic-alternant} is the exact arithmetic gap.  The real optimum can be
far below one because $S$ is large.  But the canonical scaling that makes its coefficients
integral multiplies the residual by the same $S$, leaving $D\ge1$.  Any proof along this line
must avoid paying all of $S$, or prove an extra divisibility of $D$ that is incompatible with
an independent upper bound.
\end{remark}

\subsection{A three-node example}

Take nodes $1,2,3$, basis $1,t$, and target $t^2$.  Then
\[
 A=\begin{pmatrix}1&1\\1&2\\1&3\end{pmatrix},
 \qquad
 f=\begin{pmatrix}1\\4\\9\end{pmatrix},
 \qquad
 C=(1,-2,1),
\]
so $D=2$ and $S=4$.  The best real linear approximant is
\[
 p_*(t)=4t-\frac72,
\]
with alternating residuals $(1/2,-1/2,1/2)$.  Clearing the cofactor denominator gives
\[
 4t^2-16t+14,
\]
whose node values are $(2,-2,2)$.  The small real error $1/2$ has become the nonzero integer
$2$.

For comparison, the best integer-coefficient linear approximation has uniform error one:
$p(t)=4t-3$ gives residuals $(0,-1,0)$.  Thus the evaluation lattice may beat the canonical
cofactor scaling, but it cannot cross the unit threshold on three nodes because equality at
all three would violate the zero bound.

\subsection{Smith normal form and the true arithmetic optimization}

Let $E=\{q_0,\ldots,q_N\}$ and consider the lattice $\mathcal L=A\Z^N\subset\Z^{N+1}$.
The exact arithmetic problem is
\begin{equation}
 \dist_{\ell^\infty}(f,\mathcal L).
 \label{eq:lattice-distance}
\end{equation}
The real minimax theorem replaces $\mathcal L$ by its real span.  Smith normal form describes
the quotient $\Z^{N+1}/\mathcal L$ and its congruence obstructions, while the cofactor vector
identifies the unique real dual direction.  A successful argument would combine:

\begin{enumerate}
\item an archimedean upper bound showing that $f$ lies within distance $<1$ of
      $\operatorname{span}_\R\mathcal L$;
\item a lattice theorem showing that the closest real point can be represented by an integral
      coefficient vector without increasing the error to one;
\item or, conversely, a congruence theorem forcing every nonzero residual to have absolute
      value much larger than the analytic upper bound.
\end{enumerate}

Theorem~\ref{thm:arithmetic-alternant} shows why a generic determinant estimate is not enough:
$D$ and $S$ must be studied together.

\section{$q$-Newton interpolation, Gaussian binomials, and the geometric-axis trap}
\label{sec:qnewton}

The dyadic certificates are initial pieces of a canonical integer-valued interpolation
calculus on a geometric progression.  Its formulas are exact and reveal both a large
arithmetic denominator budget and a matching extrapolation explosion.

\subsection{The normalized geometric Newton basis}

Fix integers $q\ge2$ and $r\ge1$.  For $k\ge0$, define
\begin{equation}
 B_k^{(q)}(X):=
 \prod_{i=0}^{k-1}\frac{X-q^i}{q^k-q^i},
 \qquad B_0^{(q)}=1,
 \label{eq:q-newton-basis}
\end{equation}
and
\begin{equation}
 a_k(q,r):=\prod_{i=0}^{k-1}(r-q^i),
 \qquad a_0(q,r)=1.
 \label{eq:q-newton-coeff}
\end{equation}
For $n,k\in\Nzero$,
\begin{equation}
 B_k^{(q)}(q^n)=
 \qbinom{n}{k}{q}
 :=\prod_{j=0}^{k-1}\frac{q^{n-j}-1}{q^{k-j}-1},
 \label{eq:gaussian-value}
\end{equation}
with the convention that the Gaussian binomial is zero when $k>n$.  Hence the basis is
integer-valued on the geometric set $\{q^n:n\ge0\}$.

\begin{theorem}[$q$-Newton interpolation]
\label{thm:q-newton}
Define
\begin{equation}
 P_d^{q,r}(X):=\sum_{k=0}^d a_k(q,r)B_k^{(q)}(X).
 \label{eq:q-newton-poly}
\end{equation}
Then:
\begin{enumerate}
\item $P_d^{q,r}(q^n)\in\Z$ for every $n\in\Nzero$;
\item $P_d^{q,r}(q^n)=r^n$ for $0\le n\le d$;
\item for every $n>d$,
\begin{equation}
 r^n-P_d^{q,r}(q^n)
 =\sum_{k=d+1}^n\qbinom{n}{k}{q}
   \prod_{i=0}^{k-1}(r-q^i);
 \label{eq:q-newton-remainder}
\end{equation}
\item at the first omitted geometric node,
\begin{equation}
 r^{d+1}-P_d^{q,r}(q^{d+1})
 =\prod_{i=0}^d(r-q^i).
 \label{eq:first-extrapolation}
\end{equation}
\end{enumerate}
\end{theorem}

\begin{proof}
Equation~\eqref{eq:gaussian-value} follows by factoring $q^i$ from numerator and denominator
and reversing the product.  Gaussian binomial coefficients are integers, proving the first
claim.

The Newton interpolation formula for the function $F(q^n)=r^n$ on the nodes
$1,q,q^2,\ldots$ has divided-difference coefficient
$\prod_{i<k}(r-q^i)$ in the normalized basis \eqref{eq:q-newton-basis}.  Equivalently, for
each $n$ the polynomial identity in $r$
\begin{equation}
 r^n=\sum_{k=0}^n\qbinom{n}{k}{q}
       \prod_{i=0}^{k-1}(r-q^i)
 \label{eq:q-newton-identity}
\end{equation}
holds because both sides have degree $n$ and agree at $r=1,q,\ldots,q^n$ by triangularity.
Truncating \eqref{eq:q-newton-identity} proves the second and third claims.  For $n=d+1$, the
remainder has a single term and the Gaussian coefficient is one.
\end{proof}

\subsection{Termination and the entire Newton shadow}

Let $\alpha=\log_qr$.  If $r=q^m$ with $m\in\Nzero$, then
$a_{m+1}(q,r)=0$ and all subsequent coefficients vanish.  The series terminates at
$P_m^{q,r}(X)=X^m$.  Conversely, if $r$ is not a nonnegative integral power of $q$, no
coefficient vanishes.

The denominator of $B_k^{(q)}$ is
\begin{equation}
 D_k(q):=\prod_{i=0}^{k-1}(q^k-q^i)
 =q^{k^2}\prod_{j=1}^k(1-q^{-j}).
 \label{eq:q-newton-denominator}
\end{equation}
For fixed $q,r$ with no terminating coefficient,
\begin{equation}
 \log|a_k(q,r)|=\frac{k(k-1)}2\log q+O_{q,r}(1),
 \qquad
 \log D_k(q)=k^2\log q+O_q(1).
 \label{eq:q-newton-growth}
\end{equation}
For fixed $X$, the $k$th term of \eqref{eq:q-newton-poly} is $O_{q,r,X}(q^{-k})$.
Consequently
\begin{equation}
 H_{q,r}(X):=\sum_{k=0}^\infty a_k(q,r)B_k^{(q)}(X)
 \label{eq:newton-shadow}
\end{equation}
converges locally uniformly to an entire function satisfying
$H_{q,r}(q^n)=r^n$ for every $n\ge0$.

When $\alpha\notin\Nzero$, this entire interpolation function is not $X^\alpha$: the latter
has a branch point at zero.  Thus the geometric data admit a canonical entire ``Newton
shadow'' different from the intended power function between the nodes.  This is a special
function version of the underdetermination inherent in one-axis interpolation.

For example,
\[
 H_{2,3}(3)=5.698459082460503\ldots,
 \qquad
 3^{\thetaq}=5.704522494691118\ldots.
\]
The difference is small but nonzero.  The approximation is not the issue; the shadow has no
known reason to share the second arithmetic axis.

\subsection{First extrapolation grows superexponentially}

Equation~\eqref{eq:first-extrapolation} has a particularly clear asymptotic consequence.
For fixed $q,r$ with $r\ne q^m$,
\begin{equation}
 \log\left|r^{d+1}-P_d^{q,r}(q^{d+1})\right|
 =\frac{d(d+1)}2\log q+O_{q,r}(1).
 \label{eq:extrapolation-growth}
\end{equation}
The same quadratic exponent that creates large admissible denominators also creates a large
first extrapolation residual.  Integer-valued interpolation on one geometric axis is
therefore abundant but not close to minimax extrapolation.

For $q=2,r=3$, the first omitted errors are exact integers:

\begin{table}[ht]
\centering
\small
\begin{tabular}{r@{\qquad}r@{\qquad}r}
\toprule
$d$ & $3^{d+1}-P_d^{2,3}(2^{d+1})$ & absolute value \\
\midrule
$0$ & $2$ & $2$ \\
$1$ & $2$ & $2$ \\
$2$ & $-2$ & $2$ \\
$3$ & $10$ & $10$ \\
$4$ & $-130$ & $130$ \\
$5$ & $3770$ & $3770$ \\
$6$ & $-229970$ & $229970$ \\
$7$ & $28746250$ & $28746250$ \\
$8$ & $-7272801250$ & $7272801250$ \\
$9$ & $3701855836250$ & $3701855836250$ \\
\bottomrule
\end{tabular}
\caption{The exact first-extrapolation product
$\prod_{i=0}^d(3-2^i)$.}
\label{tab:qnewton-errors}
\end{table}

\subsection{The cross-axis criterion}

Return to a hypothetical $M=2^\beta\in\Z$.  The polynomial
$P_d^{2,M}(X)$ is integer-valued on every dyadic node and agrees with $X^\beta$ at
$1,2,4,\ldots,2^d$.  The second condition $N=3^\beta\in\Z$ would be contradicted by a single
cross-axis estimate.

\begin{criterion}[Cross-axis $q$-Newton certificate]
\label{crit:cross-axis}
Write $P_d^{2,M}(3)=a/b$ in lowest terms.  If
\[
 \left|3^\beta-\frac ab\right|<\frac1b,
\]
then $3^\beta\notin\Z$.
\end{criterion}

\begin{proof}
This is Theorem~\ref{thm:dyadic-certificate}, since as a polynomial in $M$ the value
$P_d^{2,M}(3)$ is precisely the interpolant to $M^{\thetaq}$ at
$M=1,2,4,\ldots,2^d$.
\end{proof}

For $d=2$,
\[
 P_2^{2,M}(3)=\frac{M^2+3M-1}{3},
\]
which is the polynomial used in Corollary~\ref{cor:small-M}.  The open problem is to choose
nodes or a more flexible integer-valued basis so that the cross-axis error beats the reduced
denominator for every non-power $M$.

\section{Integer-valued polynomials on the $3$-smooth semigroup}
\label{sec:int-valued}

The natural arithmetic class is
\begin{equation}
 \Int(\Bsem,\Z):=
 \{P\in\Q[T]:P(q)\in\Z\text{ for every }q\in\Bsem\}.
 \label{eq:int-ring}
\end{equation}
It is much larger than $\Z[T]$ and can absorb denominators through congruences rather than
through coefficient clearing.

\subsection{Universal congruence polynomials}

For a prime $p\notin\{2,3\}$ and $k\ge1$, every $q\in\Bsem$ is a unit modulo $p^k$.
Euler's theorem gives
\[
 q^{\varphi(p^k)}\equiv1\pmod{p^k}.
\]
Hence
\begin{equation}
 U_{p,k}(T):=\frac{T^{\varphi(p^k)}-1}{p^k}
 \in\Int(\Bsem,\Z).
 \label{eq:universal-congruence}
\end{equation}
This elementary family proves that coefficient denominators need not be paid by a factor
$X^d$ at the evaluation stage.  Sparse-set congruences can pay them uniformly.

The efficiency of this particular family is limited:
\[
 \frac{\log p^k}{\deg U_{p,k}}
 =\frac{k\log p}{p^{k-1}(p-1)}.
\]
It decays rapidly in $p$ and $k$.  Products over many primes create larger denominators but
also larger degrees.  The relevant extremal question is global.

\subsection{Bhargava orderings and generalized factorials}

Bhargava's $p$-orderings attach to an arbitrary subset $S$ of a Dedekind domain a generalized
factorial sequence measuring the largest denominators compatible with integer-valuedness
\cite{Bhargava1997}.  Recent $B$-ordering work supplies simultaneous variants for collections
of ideals \cite{LagariasYangjit2025}.  For $S=\Bsem$, the desired invariant is the
characteristic fractional ideal
\[
 \mathfrak I_d(S)=
 \bigl\{a_d:\exists\,P(T)=\sum_{j=0}^d a_jT^j\in\Int(S,\Z)\bigr\}.
\]
Equivalently, one wants a regular basis
$F_0,F_1,\ldots$ of $\Int(S,\Z)$ with controlled denominators, heights, and norms on
multiplicative blocks.

The geometric Newton basis demonstrates the possible scale.  On
$S_q=\{q^n:n\ge0\}$, the degree-$k$ denominator
$D_k(q)$ in \eqref{eq:q-newton-denominator} satisfies
\[
 \log D_k(q)=k^2\log q+O(1).
\]
Thus a sparse set can support quadratic-in-degree denominator growth.  Since the naive
normalization tax has logarithm $d\log X$ and the interesting regime has $d\asymp\log X$,
both live on the $(\log X)^2$ scale.  This is the first reason the integer-valued-polynomial
route is not ruled out by Proposition~\ref{prop:clearing-tax}.

The second reason is also a warning: for the canonical $q$-Newton interpolant, the Newton
coefficient consumes about half of that quadratic budget, and the first extrapolation error
grows on the same scale.  The full semigroup $\Bsem$ may provide better mixed-base
congruences, but they must be optimized jointly with archimedean norm.

\subsection{The exact arithmetic Bernstein criterion}

The open construction can be stated as a theorem with a single missing hypothesis.

\begin{criterion}[Arithmetic Bernstein criterion]
\label{crit:arithmetic-bernstein}
Let $\beta\notin\Nzero$ satisfy $q^\beta\in\Z$ for all $q\in\Bsem$.  Suppose that for some
fixed $c>1$ there are arbitrarily large $X$, degrees $d=d(X)$, and polynomials
$P_X\in\Int(\Bsem,\Z)$ such that
\begin{enumerate}
\item $\deg P_X\le d$;
\item $\#(\Bsem\cap[X,cX])\ge d+2$;
\item
\[
 \max_{q\in\Bsem\cap[X,cX]}|q^\beta-P_X(q)|<1.
\]
\end{enumerate}
Then no such $\beta$ exists.
\end{criterion}

\begin{proof}
Apply Theorem~\ref{thm:locking} to the node set
$\Bsem\cap[X,cX]$.
\end{proof}

For $M=3$, Theorem~\ref{thm:optimal-block} says the degree and node counts leave a positive
linear-in-$\log X$ margin at $c=c_*$.  Thus it would be enough to prove an
integer-valued analogue of Bernstein approximation with no loss in the exponential base.
For $M\ge5$, the same function space cannot work even under ideal arithmetic
integerization; one needs rational or M\"untz spaces, additional nodes, or a nonarchimedean
amplifier.

\subsection{A concrete $p$-ordering computation program}

For each prime $p$, a $p$-ordering of $\Bsem$ greedily chooses
$a_0,a_1,\ldots\in\Bsem$ so that
\[
 v_p\!\left(\prod_{i<n}(a_n-a_i)\right)
\]
is minimal at each stage.  The resulting valuation is independent of the choices and gives
the $p$-part of the generalized factorial.  A finite rigorous program can:

\begin{enumerate}
\item enumerate $3$-smooth numbers up to a certified cutoff;
\item compute greedy minima for each $p$ up to a degree $d$;
\item use periodicity of $2^a3^b\bmod p^k$ to prove that the finite search is complete;
\item combine the local factorials into a candidate regular basis;
\item solve a linear or lattice optimization for minimal uniform error on
      $\Bsem\cap[X,cX]$.
\end{enumerate}

The key theorem to seek is not merely a formula for the generalized factorial.  It is a
simultaneous bound
\begin{equation}
 \inf_{P\in\Int(\Bsem,\Z),\,\deg P\le d}
 \max_{q\in\Bsem\cap[X,cX]}|q^\beta-P(q)|
 \le \exp\bigl(\beta\log X-d\log\rho(c)+o(\log X)\bigr),
 \label{eq:int-bernstein-target}
\end{equation}
with a loss smaller than the positive margin in Table~\ref{tab:efficiency}.  No such theorem
is presently known.

\subsection{Why regular bases and $p$-adic Fourier theory matter}

Berger and Sprang show that regular bases of rings of integer-valued polynomials can be
encoded by $p$-adic Fourier expansions and used to derive numerical criteria for arithmetic
properties \cite{BergerSprang2025}.  In the present setting, a regular basis adapted to
$\Bsem$ could do two things at once:

\begin{itemize}
\item make integer-valuedness automatic, replacing a hard global congruence constraint by
      integral basis coefficients;
\item expose the archimedean norm of each basis element on $[X,cX]$, turning
      \eqref{eq:int-bernstein-target} into a weighted coefficient approximation problem.
\end{itemize}

The obstacle is that $\Bsem$ is sparse in the real topology but highly recurrent in every
$p$-adic unit group away from $2$ and $3$.  That mixture is exactly what generalized
factorials are designed to encode.

\section{Rational approximation and the moving denominator threshold}
\label{sec:rational}

Polynomial approximation is only one member of the classical hierarchy.  Rational functions
can place poles near the branch cut of $t^\beta$ and may converge substantially faster on a
fixed compact interval.  In the arithmetic problem, however, the comparison scale is no
longer the fixed number one.  The denominator of the approximant is evaluated at the
arithmetic node, so the locking threshold moves with the node.

\subsection{An exact rational locking theorem}

\begin{theorem}[Rational-function locking]
\label{thm:rational-locking}
Let $\alpha>0$ be nonintegral and let $E\subset\N$ be a finite set such that
$q^\alpha\in\Z$ for every $q\in E$.  Let $P,Q\in\Z[T]$ be nonzero with
\[
 \deg P\le m,\qquad \deg Q\le n,
\]
and suppose $Q(q)\ne0$ for all $q\in E$.  If
\begin{equation}
 \left|q^\alpha-\frac{P(q)}{Q(q)}\right|<\frac1{|Q(q)|}
 \qquad(q\in E),
 \label{eq:rational-lock}
\end{equation}
then $|E|\le m+n+1$.  Consequently, finding $m+n+2$ such nodes satisfying
\eqref{eq:rational-lock} excludes the assumed integrality.
\end{theorem}

\begin{proof}
For every $q\in E$, the quantity
\[
 q^\alpha Q(q)-P(q)
\]
is an integer.  Inequality~\eqref{eq:rational-lock} forces that integer to vanish.  Hence the
generalized polynomial
\[
 F(t)=t^\alpha Q(t)-P(t)
\]
has $|E|$ positive zeros.  Its exponents lie in the disjoint sets
$\{\alpha,\alpha+1,\ldots,\alpha+n\}$ and $\{0,1,\ldots,m\}$, which are disjoint because
$\alpha\notin\Z$.  Thus $F$ has at most $m+n+2$ nonzero terms and, by the generalized
Descartes--Chebyshev zero theorem, at most $m+n+1$ positive zeros.  The function cannot vanish
identically because the two exponent sets are disjoint.
\end{proof}

For the conjecture one may take $\alpha=\beta$ and $E\subset\Bsem$, or
$\alpha=\thetaq$ and $E$ in the conditional candidate semigroup
$\{2^aM^b:a,b\ge0\}$.  The theorem is exact and does not assume uniform approximation on the
whole interval.

\subsection{Why ordinary rational convergence is not the right metric}

On a fixed interval $[1,c]$, a type-$(m,n)$ rational approximant may have a very small real
error.  After transporting it to $[X,cX]$, however, an integer polynomial denominator of
degree $n$ is typically of size $X^n$ at the nodes.  The arithmetic requirement is therefore
not merely
\[
 \|t^\beta-R(t)\|=o(1),
\]
but a weighted estimate
\begin{equation}
 \max_{q\in E}|Q(q)|\,
 \left|q^\beta-\frac{P(q)}{Q(q)}\right|<1.
 \label{eq:weighted-rational-target}
\end{equation}
If $n\asymp\log X$ and $|Q(q)|$ has its generic $X^n$ scale, the required error is
$\exp(-\Theta((\log X)^2))$.  Fixed-block polynomial approximation supplies only
$\exp(-\Theta(n))$, and even many standard rational schemes supply an exponential in the
degree rather than the required exponential in $n\log X$.

This does not prove that every rational scheme fails.  It identifies the missing property:
the denominator must be arithmetically tiny on the selected smooth nodes while the numerator
still approximates $t^\beta Q(t)$.  In other words, the problem is a weighted rational
approximation problem with a denominator constrained simultaneously in the real and
nonarchimedean metrics.

\begin{criterion}[Weighted rational certificate]
\label{crit:weighted-rational}
A counterexample is impossible if, for arbitrarily large $X$, there are
$P_X,Q_X\in\Z[T]$ and a set
$E_X\subset\Bsem\cap[X,cX]$ such that
\begin{enumerate}
\item $|E_X|>\deg P_X+\deg Q_X+1$;
\item $Q_X$ has no zero on $E_X$;
\item the weighted bound~\eqref{eq:weighted-rational-target} holds at every point of $E_X$.
\end{enumerate}
\end{criterion}

\subsection{Pad\'e approximants and the coefficient-field obstruction}

The local Pad\'e approximants to $(1+z)^\beta$ have coefficients rational in $\beta$.
For algebraic $\beta$ that feature is useful, and for rational $\beta$ it gives classical
irrationality constructions.  Here a hypothetical counterexample has transcendental
$\beta$.  The Pad\'e coefficients therefore lie in $\Q(\beta)$ rather than $\Q$, and their
evaluation at an integer node is not an integer or even an algebraic number with controlled
denominator.  Clearing the coefficient field would require a polynomial relation for
$\beta$, precisely what is unavailable.

Approximating $t^{\thetaq}$ avoids an unknown exponent but replaces the coefficient field by
$\Q(\thetaq)$.  Since $\thetaq$ is transcendental, the same obstruction remains.  A successful
Pad\'e construction must therefore be adapted to the exponential identities
$2^{\thetaq}=3$ and $M^{\thetaq}=N$, not merely to the Taylor series of a generic power.

\begin{statusbox}{Theorem~\ref{thm:rational-locking} and
Criterion~\ref{crit:weighted-rational} are \statusnew.  The diagnosis of standard Pad\'e
coefficients is structural: it says why the usual algebraic-exponent machinery does not
specialize, not that every possible rational auxiliary function is excluded.}
\end{statusbox}


\section{Hermite--Pad\'e auxiliary functions and the four-exponentials wall}
\label{sec:hermite-pade}

Recent work of Fischler and Rivoal obtains a new transcendence measure for values of the
exponential function at nonzero algebraic arguments by combining Hermite--Pad\'e approximants
to powers of the exponential function with a carefully optimized Siegel determinant
\cite{FischlerRivoal2025}.  Faverjon and Adamczewski likewise obtain explicit algebraic
independence measures for values of $E$-functions and Mahler functions at algebraic points
\cite{AdamczewskiFaverjon2025}.  These are major improvements in quantitative auxiliary
function technology.  The present problem sits just outside their arithmetic hypotheses.

\subsection{The algebraic-point mismatch}

The four relevant exponentials can be written as
\[
 e^{\log2}=2,\qquad e^{\log3}=3,
 \qquad e^{\beta\log2}=M,
 \qquad e^{\beta\log3}=N.
\]
All four values are algebraic, but the arguments $\log2,\log3,\beta\log2,\beta\log3$ are not
algebraic.  Hermite--Pad\'e theorems for $E$-functions obtain arithmetic control by evaluating
a differential system with algebraic coefficients at an algebraic point.  Here the natural
differential systems
\[
 y'=\log2\,y,
 \qquad
 y'=\log3\,y
\]
already have transcendental coefficients, and the common evaluation point $z=\beta$ is itself
transcendental under a counterexample.  Neither side of the usual arithmetic input survives.

This is not a technical inconvenience.  Replacing the logarithms by formal algebraic
parameters and later specializing them would require proving that a nonzero formal auxiliary
determinant cannot vanish at $(\log2,\log3)$.  That is a statement about products of
logarithms of algebraic numbers, exactly the structural-rank wall isolated in the supplied
report and in the Matrix Coefficient program of Dasgupta--Kakde
\cite{DasguptaKakde2024,DasguptaKakde2026}.

\subsection{Exponential polynomials with integral values}

A direct auxiliary function has the form
\begin{equation}
 F(z)=\sum_{(a,b)\in\Omega}c_{a,b}
       e^{(a\log2+b\log3)z}
     =\sum_{(a,b)\in\Omega}c_{a,b}2^{az}3^{bz},
 \qquad c_{a,b}\in\Z.
 \label{eq:direct-exp-poly}
\end{equation}
At $z=\beta$ every term is an integer, so $F(\beta)\in\Z$.  If one could arrange
$0<|F(\beta)|<1$, the conjecture would follow.

High-order vanishing at zero would impose
\begin{equation}
 \sum_{(a,b)\in\Omega}c_{a,b}(a\log2+b\log3)^j=0
 \qquad(0\le j<L).
 \label{eq:moment-log}
\end{equation}
For algebraic frequencies, Siegel's lemma can solve analogous systems over a number field.
Here the coefficients of~\eqref{eq:moment-log} are polynomials in $\log2$ and $\log3$.
Solving them with integer $c_{a,b}$ without treating those logarithms as independent is
already a logarithmic-algebraic problem.  Treating them as formally independent replaces each
condition by all its binomial coefficient conditions
\[
 \sum_{(a,b)\in\Omega}c_{a,b}a^rb^{j-r}=0,
 \qquad 0\le r\le j,
\]
and consumes about $L^2/2$ degrees of freedom.  This quadratic cost is the auxiliary-function
shadow of the critical $2\times2$ four-exponentials geometry.

\subsection{The polynomial-on-a-curve version}

Writing $X=2^z$ and $Y=3^z$, equation~\eqref{eq:direct-exp-poly} becomes
$F(z)=Q(X,Y)$ with $Q\in\Z[X,Y]$.  A counterexample supplies the integral curve point
$(M,N)=(2^\beta,3^\beta)$ on
\[
 Y=X^{\thetaq}.
\]
The inherited sparse-interpolation theorem says that a nonzero polynomial with $s$ monomials
has at most $s-1$ distinct positive zeros on an irrational power curve.  Thus interpolation
through $R$ positive nodes requires at least $R+1$ monomials.  The supplied report already
identifies the exact one-monomial saving that would prove the conjecture.  Hermite--Pad\'e
language does not erase that deficit; it repackages it as the need for one more exact moment
condition than the available support supplies.

\subsection{What the recent transcendence papers do supply}

The recent literature is nevertheless instructive in three ways.

\begin{enumerate}
\item Fischler--Rivoal show that accessory parameters in a Siegel determinant can improve a
      transcendence measure beyond the previously expected barrier.  This encourages a
      search for an accessory parameter adapted to the two prime valuations rather than to a
      single Archimedean norm.
\item Faverjon--Adamczewski show how functional equations and differential systems can turn
      algebraic independence into explicit Liouville-type inequalities.  A comparable
      inequality for the specialized logarithmic matrix would immediately be relevant, but
      no such theorem is presently known.
\item Massold's 2025 work develops lower bounds for transcendence degrees of fields generated
      by exponentials of products and reduces stronger estimates to a torus-intersection
      problem \cite{Massold2025}.  The $2\times2$ specialization remains on the conjectural
      side; it does not furnish the missing auxiliary function here.
\end{enumerate}

\begin{statusbox}{This section is an audited transfer analysis, not a negative theorem about
all Hermite--Pad\'e methods.  The exact obstruction is that the natural arguments and
coefficient matrices are transcendental, while current quantitative $E$-function theorems
require algebraic arithmetic data.}
\end{statusbox}


\section{Orthogonal polynomials, Christoffel functions, and lattice rounding}
\label{sec:orthogonal}

Uniform approximation is the right norm for the one-unit locking principle, but least-squares
approximation has a useful exact determinant formula.  It exposes another duality between
approximation and the integer interpolation determinants of the supplied report.

\subsection{A Gram-determinant identity}

Let $E=\{q_0,\ldots,q_N\}$, let
$A=(\phi_j(q_i))_{0\le i\le N,\,0\le j<N}$, and let
$f=(\psi(q_i))_{0\le i\le N}$.  Give the nodes positive weights
$w_i$ and write $W=\operatorname{diag}(w_0,\ldots,w_N)$ and $V=[A\ f]$.

\begin{proposition}[Least-squares determinant quotient]
\label{prop:gram-quotient}
If $A$ has full column rank, then
\begin{equation}
 \min_{c\in\R^N}(f-Ac)^TW(f-Ac)
 =\frac{\det(V^TWV)}{\det(A^TWA)}.
 \label{eq:gram-quotient}
\end{equation}
\end{proposition}

\begin{proof}
The normal equations give the minimizer
$c=(A^TWA)^{-1}A^TWf$.  The minimum is the Schur complement
\[
 f^TWf-f^TWA(A^TWA)^{-1}A^TWf
\]
of $A^TWA$ in $V^TWV$.  The determinant formula for a Schur complement gives
\eqref{eq:gram-quotient}.
\end{proof}

For polynomial bases, the denominator is a moment or Gram determinant and its reciprocal is a
Christoffel-type quantity.  Choosing weights to minimize the largest leverage score is the
least-squares analogue of choosing Fekete or alternation nodes in the uniform norm.

\subsection{The arithmetic gap in quadratic norm}

Under a counterexample, choose $q_i\in\Bsem$ and exponents from
$\Nzero+\Nzero\beta$.  Then $A$ and $f$ have integer entries.  If the weights are integers,
both determinants in~\eqref{eq:gram-quotient} are integers.  This looks promising, but the
quotient can be much smaller than one because the denominator determinant may be enormous.
Moreover, the real minimizer generally has rational rather than integral coefficients.
Clearing its denominator again multiplies the residual to a nonzero integer vector.

The correct arithmetic object is the evaluation lattice
\begin{equation}
 \cL_E=A\Z^N\subset\Z^{N+1}.
 \label{eq:evaluation-lattice}
\end{equation}
The desired approximation asks for the distance of the integral target $f$ from this lattice,
not merely from the real subspace $A\R^N$:
\begin{equation}
 \delta_\infty(E,A;f)
 :=\min_{z\in\Z^N}\|f-Az\|_\infty.
 \label{eq:evaluation-lattice-distance}
\end{equation}
When $|E|=N+1$, Theorem~\ref{thm:arithmetic-alternant} supplies the primitive dual
functional.  Let $C=(C_0,\ldots,C_N)$ be its cofactor vector and let
$g=\gcd(C_0,\ldots,C_N)$.  Then $C/g$ is primitive and every residual $r=f-Az$ satisfies
\begin{equation}
 \frac{C}{g}\cdot r=\frac{D}{g}.
 \label{eq:primitive-functional}
\end{equation}
The real lower bound is $|D|/S$; the integral lower bound is the integer optimization problem
of minimizing $\|r\|_\infty$ subject to~\eqref{eq:primitive-functional} and the torsion
congruences of the Smith normal form of $A$.

\begin{criterion}[Smith-coset criterion]
\label{crit:smith-coset}
A polynomial-approximation proof would follow if, for some counterexample candidate and some
node/basis family with more nodes than basis functions, one could prove from the Smith normal
form that the coset $f+\cL_E$ contains the zero vector while the Chebyshev zero theorem proves
that it cannot.  Equivalently, a contradiction arises if the arithmetic congruences force an
exact interpolant of insufficient dimension.
\end{criterion}

This criterion deliberately asks for exact membership, not a small real error.  Since all
residual coordinates are integers, every nonzero coset representative has
$\ell^\infty$ norm at least one.

\subsection{Why independent coefficient rounding is too noisy}

Let $c^*$ be a real minimax or least-squares coefficient vector and round each coordinate
independently to an integer with mean $c_j^*$.  At node $q_i$ the rounding perturbation is
\[
 Z_i=\sum_j A_{ij}\xi_j,
 \qquad |\xi_j|\le1.
\]
Even under unbiased Bernoulli rounding,
\[
 \operatorname{Var}(Z_i)
 =\sum_j A_{ij}^2\operatorname{Var}(\xi_j)
 \le\frac14\sum_jA_{ij}^2.
\]
For monomial or M\"untz bases on nodes of size $X$, the right side is dominated by powers of
$X$ and is vastly larger than one.  Orthogonalizing the columns reduces the real condition
number, but the orthogonal basis has nonintegral coefficients and destroys the evaluation
lattice.  Thus generic randomized rounding attacks the wrong lattice.

A potentially relevant replacement is discrepancy minimization in a basis already adapted to
$\Int(\Bsem,\Z)$: one would round coefficients in a regular integer-valued basis whose
vectors are small on the selected block.  This is exactly where generalized factorials and
$p$-adic Fourier bases can matter.


\section{A M\"untz-space formulation and two exact approximation floors}
\label{sec:muntz}

The conditional exponent semigroup
\[
 \Lambda_\beta=\{n+k\beta:n,k\in\Nzero\}
\]
provides far more functions than the ordinary powers.  Every finite collection of distinct
exponents in $\Lambda_\beta$ is a M\"untz--Chebyshev system on $(0,\infty)$, and every such
monomial is integral on the base semigroup $\Bsem$ under a counterexample.

\subsection{The arithmetic M\"untz lattice}

For finite $\Lambda\subset\Lambda_\beta$, define
\[
 \mathcal M_\Z(\Lambda)
 =\left\{\sum_{\lambda\in\Lambda}c_\lambda t^\lambda:
               c_\lambda\in\Z\right\}.
\]
If $q=2^a3^b$, then
\[
 q^{n+k\beta}=q^n(M^aN^b)^k\in\Z.
\]
Hence every $F\in\mathcal M_\Z(\Lambda)$ is integer-valued on $\Bsem$.

\begin{theorem}[Unit floor for integer M\"untz polynomials]
\label{thm:muntz-unit-floor}
Let $\Lambda$ contain $s$ distinct exponents and let
$E\subset\Bsem$ have at least $s$ points.  Under a nonintegral counterexample,
\[
 0\ne F\in\mathcal M_\Z(\Lambda)
 \quad\Longrightarrow\quad
 \max_{q\in E}|F(q)|\ge1.
\]
More generally, if $\mu\in\Lambda_\beta\setminus\Lambda$ and $|E|\ge s+1$, then
\begin{equation}
 \inf_{(c_\lambda)\in\Z^\Lambda}
 \max_{q\in E}
 \left|q^\mu-\sum_{\lambda\in\Lambda}c_\lambda q^\lambda\right|
 \ge1.
 \label{eq:muntz-inhom-floor}
\end{equation}
\end{theorem}

\begin{proof}
All displayed values are integers.  A strict bound below one would therefore force a zero at
every node.  A nonzero $s$-term M\"untz polynomial has at most $s-1$ positive zeros, while the
inhomogeneous residual has at most $s+1$ terms and at most $s$ positive zeros.
\end{proof}

This theorem is simple but conceptually decisive.  Classical M\"untz approximation permits
arbitrarily small real errors in many settings; the arithmetic coefficient lattice has a hard
floor exactly at one.  The problem is to build a basis in which real approximation and
integer-valuedness coexist before coefficient rounding.

\subsection{A continuous-versus-arithmetic dichotomy}

Because ordinary nonnegative integers already lie in $\Lambda_\beta$, the real span contains
all algebraic polynomials and is dense on every compact interval.  Thus lack of analytic
flexibility is not the issue.  The obstruction is entirely in the coefficient lattice.
One may summarize it as
\[
 \underbrace{\operatorname{dist}_{\infty}
  (t^\mu,\operatorname{span}_{\R}\{t^\lambda\})}_{\text{can be very small}}
 \qquad\text{versus}\qquad
 \underbrace{\operatorname{dist}_{\infty}
  (t^\mu,\mathcal M_\Z(\Lambda))}_{\text{at least }1}.
\]
The cofactor formula of Section~\ref{sec:minimax} computes the first quantity on a finite
alternation set; the Smith form of the evaluation matrix controls the second.

\subsection{Signed exponents and the denominator price}

Allowing negative exponents would create frequencies close to $\beta$ through good rational
approximations $n+k\beta\approx0$ and could improve real approximation dramatically.  But
$q^{n+k\beta}$ then acquires denominators on $\Bsem$.  Clearing them reintroduces prime powers
whose logarithm is proportional to the size of the negative exponent.  This is the M\"untz
version of the normalization tax: frequency localization and arithmetic integrality pull in
opposite directions.

A promising restricted question is to allow Laurent exponents while keeping all denominators
supported only at $2$ and $3$, then seek a compensating $2$-adic and $3$-adic divisor in the
cofactor determinant.  The supplied report's signed-progression and valuation modules provide
part of the bookkeeping, but no leading-order cancellation is currently known.


\section{Computational reconnaissance and reproducible certificates}
\label{sec:computation}

The computations in this continuation are diagnostic rather than evidence for a full proof.
They are small enough to reproduce with exact integer or rational arithmetic, except for the
single displayed value of the entire Newton shadow.

\subsection{Counting $3$-smooth nodes in the optimal block}

The following table enumerates $\Bsem\cap[X,c_*X]$ exactly and compares it with the leading
term from Proposition~\ref{prop:thin-strip}.  The modest oscillation around the asymptotic is
visible even at moderate scales.

\begin{table}[ht]
\centering
\small
\begin{tabular}{r@{\qquad}r@{\qquad}r@{\qquad}r}
\toprule
$X$ & exact count & leading term & difference \\
\midrule
$10^2$  & $15$ & $10.6602$ & $4.3398$ \\
$10^3$  & $19$ & $15.9903$ & $3.0097$ \\
$10^4$  & $25$ & $21.3204$ & $3.6796$ \\
$10^6$  & $36$ & $31.9806$ & $4.0194$ \\
$10^8$  & $48$ & $42.6408$ & $5.3592$ \\
$10^{10}$ & $57$ & $53.3011$ & $3.6989$ \\
$10^{12}$ & $68$ & $63.9613$ & $4.0387$ \\
\bottomrule
\end{tabular}
\caption{Exact $3$-smooth node counts for $c_*=3+2\sqrt2$.}
\label{tab:node-counts}
\end{table}

The asymptotic coefficient $2.3148353905\ldots$ is large enough to beat the real polynomial
degree coefficient $1.7982868151\ldots$ only in the $M=3$ branch.  Table~\ref{tab:node-counts}
does not address integer-valuedness.

\subsection{Exact certificate generator}

The following compact script constructs a rational dyadic interpolant, evaluates it at a
candidate $M$, and reports the reduced denominator.  The analytic interval needed by
Theorem~\ref{thm:dyadic-certificate} should be proved separately by integer power comparisons.

\begin{lstlisting}[style=pythonstyle,caption={Exact dyadic interpolation certificate generator.}]
from fractions import Fraction


def interpolate(nodes, values):
    """Return coefficients low-to-high of the exact interpolant."""
    p = [Fraction(0)]
    for i, xi in enumerate(nodes):
        basis = [Fraction(1)]
        denom = Fraction(1)
        for j, xj in enumerate(nodes):
            if i == j:
                continue
            out = [Fraction(0)] * (len(basis) + 1)
            for k, a in enumerate(basis):
                out[k] -= xj * a
                out[k + 1] += a
            basis = out
            denom *= xi - xj
        scale = Fraction(values[i], 1) / denom
        if len(p) < len(basis):
            p += [Fraction(0)] * (len(basis) - len(p))
        for k, a in enumerate(basis):
            p[k] += scale * a
    return p


def evaluate(coeffs, x):
    y = Fraction(0)
    for a in reversed(coeffs):
        y = y * x + a
    return y


exponents = [0, 1, 2]
nodes = [2**n for n in exponents]
values = [3**n for n in exponents]
poly = interpolate(nodes, values)
for M in [3, 5, 6, 7]:
    value = evaluate(poly, M)
    print(M, value, value.denominator)
\end{lstlisting}

\subsection{A finite $p$-ordering experiment}

For a finite set $S_H=\{2^a3^b:0\le a,b\le H\}$ and a prime $p$, one can greedily construct
a $p$-ordering by selecting the next element $x$ minimizing
\[
 v_p\!\left(\prod_{y\text{ already selected}}(x-y)\right).
\]
Repeating this for many primes produces an empirical generalized-factorial profile.  The
experiment should record:

\begin{enumerate}
\item the greedy sequence and its valuation increments;
\item stabilization as $H$ grows;
\item the leading coefficient denominator permitted in degree $d$;
\item the sup norm on $\Bsem\cap[X,c_*X]$ of the corresponding regular-basis element;
\item the ratio between that norm and the Bernstein minimax error.
\end{enumerate}

The key output is not merely a large denominator.  It is a basis element having a large
admissible denominator and a small block norm at the same time.  Geometric $q$-Newton bases
supply the first property but fail the second by
Theorem~\ref{thm:q-newton}; ordinary Chebyshev polynomials supply the second but fail the first.

\subsection{Numerical trust boundary}

All tables except the Newton-shadow value are derived from exact integers or rational
arithmetic and rounded only for presentation.  No numerical calculation in this continuation
is used as a premise in a proof of the conjecture.  The symbolic exclusions in
Corollary~\ref{cor:small-M} reduce entirely to the displayed integer inequalities.


\section{Exact sufficient conditions extracted from approximation theory}
\label{sec:criteria}

The preceding analysis yields several precise statements whose proof would settle the
conjecture or a substantial branch of it.  They are listed here in decreasing order of
immediate plausibility.

\subsection{The generator-three integer Bernstein target}

\begin{criterion}[Optimal-block integer Bernstein target]
\label{crit:M3-bernstein}
Let $c_*=3+2\sqrt2$.  If there are arbitrarily large $X$ and polynomials
$P_X\in\Int(\Bsem,\Z)$ such that
\[
 \deg P_X<\#(\Bsem\cap[X,c_*X])-1
\]
and
\[
 \max_{q\in\Bsem\cap[X,c_*X]}|q^{\thetaq}-P_X(q)|<1,
\]
then $3^{\thetaq}\notin\Z$.
\end{criterion}


\begin{proof}
Apply Theorem~\ref{thm:locking} with $\alpha=\thetaq$ and the indicated node set.
\end{proof}

This criterion focuses on the only nonintegral generator for which ordinary Bernstein
geometry has enough asymptotic nodes.  It is therefore the natural first branch for a
$p$-ordering or integer-Chebyshev attack.

\subsection{A generalized-factorial norm target}

Let $F_d$ be a degree-$d$ regular-basis element of $\Int(\Bsem,\Z)$, normalized by its leading
coefficient.  A sufficient quantitative theorem would be a bound of the form
\begin{equation}
 \|F_d\|_{[X,c_*X]}
 \le \exp(o(d))\,\rho(c_*)^{-d}X^d
 \label{eq:factorial-norm-target}
\end{equation}
together with enough control of the lower-degree basis elements to approximate
$X^{\thetaq}t^{\thetaq}$.  The $X^d$ factor in~\eqref{eq:factorial-norm-target} is unavoidable
for a monic polynomial; the generalized factorial must cancel it in the leading coefficient.
Thus one needs
\[
 \log d!_{\Bsem}\sim d\log X
\]
on the chosen degree/scale relation, with the same basis retaining the Chebyshev norm profile.
This is an adelic capacity problem rather than a purely real one.

\subsection{A cofactor-denominator saving}

In the notation of Theorem~\ref{thm:cofactor-minimax}, the best real residual is $D/S$ and the
cleared arithmetic residual is $D$.  A proof would follow from a construction in which the
minimax coefficient vector is already integral, or more generally in which its denominator
is strictly smaller than the natural cofactor sum and the remaining residual is subunit on
one additional arithmetic node.

\begin{criterion}[Cofactor saving]
\label{crit:cofactor-saving}
Suppose a family of smooth nodes and integral M\"untz basis functions admits a best
alternating residual whose coefficients all have a common denominator $H<S$ and such that
$H(D/S)<1$.  If the corresponding coefficient clearing preserves target coefficient one,
then no counterexample exists.
\end{criterion}

The last clause is essential: multiplying the entire residual by $H$ changes the target from
$t^\mu$ to $Ht^\mu$.  One needs a congruence that clears the approximating coefficients while
leaving the target coefficient congruent to one.

\subsection{A cross-axis Newton target}

For each non-power integer $M$, choose dyadic nodes
$2^{n_0},\ldots,2^{n_d}$ and let $I_{\mathbf n}$ interpolate
$(2^{n_j},3^{n_j})$.  Theorem~\ref{thm:dyadic-certificate} reduces the conjecture to finding a
node set for which
\[
 \left|M^{\thetaq}-I_{\mathbf n}(M)\right|
 <\frac1{\den(I_{\mathbf n}(M))}.
\]
The node set may depend on $M$.  The standard consecutive dyadic nodes have excellent
integer-valuedness but poor extrapolation.  Nonconsecutive nodes, Leja-type node selection,
or a mixture of dyadic and triadic axes are concrete alternatives.

\subsection{A weighted rational target}

Criterion~\ref{crit:weighted-rational} asks for rational functions whose denominator is small
at smooth nodes.  A natural sufficient theorem would construct $Q$ from cyclotomic or
universal-congruence factors so that $Q(q)$ has controlled prime divisibility for all
$q\in\Bsem$, while $P/Q$ is a near-best rational approximation to $t^\beta$.  Such a theorem
would combine potential theory with the generalized factorial of $\Bsem$.

\subsection{A new logarithmic auxiliary-function theorem}

Finally, any theorem producing an integer exponential polynomial
$F$ of the form~\eqref{eq:direct-exp-poly} with
$0<|F(\beta)|<1$ would solve the problem immediately.  The moment analysis shows that the needed
coefficient estimate must exploit the specialization $(\log2,\log3)$, not treat the two logs
as generic independent variables.  This is the point at which approximation theory meets the
Matrix Coefficient Conjecture.


\section{Formalization blueprint}
\label{sec:formalization}

The supplied repository already contains APIs for exponential-polynomial zero bounds,
integer determinants, finite candidates, and several semigroup grids.  The continuation can
be formalized in layers.

\subsection{Low-cost exact modules}

\begin{description}[style=nextline]
\item[\texttt{OptimalBernsteinBlock}]
Define $\rho(c)$ and $\Psi(c)$; prove the identity $\sinh y\sinh u=1$, the log-convexity
inequality, and the maximizer $c_*=3+2\sqrt2$.  This is elementary real analysis once the
convexity lemma is available.

\item[\texttt{DyadicInterpolationCertificate}]
Formalize rational Lagrange interpolation on dyadic nodes, the Chebyshev zero-count endpoint,
and the denominator locking lemma.  Instantiate it for $M=3,5,6,7$ using norm-num-certified
integer powers.

\item[\texttt{DiscreteCofactorAlternation}]
Formalize the cofactor vector, $A^TC=0$, the $\ell^1$ dual formula, and
$B(z,D)^T=Sf$.  The arithmetic alternant is pure linear algebra over $\Z$ and $\Q$.

\item[\texttt{GeometricNewtonBasis}]
Define $B_k^{(q)}$, prove evaluation as a Gaussian binomial coefficient, derive the Newton
coefficients $a_k$, and prove the first-extrapolation product.

\item[\texttt{RationalPowerLocking}]
Formalize Theorem~\ref{thm:rational-locking} using the existing exponential-polynomial zero
API.
\end{description}

\subsection{Medium-cost analytic modules}

\begin{description}[style=nextline]
\item[\texttt{PowerBernsteinRate}]
Formalize the affine map to $[-1,1]$, Bernstein ellipses, the branch obstruction at zero, and
the root asymptotic for best approximation.  This likely requires a significant complex
approximation library.

\item[\texttt{SmoothThinStripCount}]
Formalize the fixed-width count of $2^a3^b$ in a multiplicative block using irrational
rotation and discrepancy.  Much of the equidistribution infrastructure already exists in the
repository.

\item[\texttt{EntireNewtonShadow}]
Prove compact convergence of the geometric Newton series and its interpolation property.
Only geometric estimates and standard uniform convergence are needed.
\end{description}

\subsection{Research-level arithmetic modules}

\begin{description}[style=nextline]
\item[\texttt{SmoothSemigroupPOrdering}]
Compute and certify local $p$-orderings of $\Bsem$, generalized factorials, and regular-basis
leading coefficients.

\item[\texttt{AdelicBlockNorm}]
Relate the generalized factorial to the real block norm of regular-basis elements.  This is
the central open approximation target.

\item[\texttt{EvaluationLatticeSmith}]
Compute Smith forms for selected smooth-node M\"untz matrices and extract exact coset
obstructions.  Finite instances are certifiable even before an asymptotic theorem is found.
\end{description}

The recommended order is exact algebra first, then finite certified experiments, then the
complex approximation asymptotic.  None of the open arithmetic modules should be represented
as a theorem until its denominator and norm estimates are both proved.


\section{Comparison of approximation approaches}
\label{sec:comparison}

{\small
\begin{longtable}{>{\raggedright\arraybackslash}p{0.17\textwidth}
                  >{\raggedright\arraybackslash}p{0.20\textwidth}
                  >{\raggedright\arraybackslash}p{0.22\textwidth}
                  >{\raggedright\arraybackslash}p{0.25\textwidth}}
\toprule
Method & Archimedean strength & Arithmetic advantage & Exact stopping point \\
\midrule
\endfirsthead
\toprule
Method & Archimedean strength & Arithmetic advantage & Exact stopping point \\
\midrule
\endhead
Bernstein polynomial & Exponential in degree with exact rate $\rho(c)^{-d}$ & None for the real minimax coefficients & Node geometry reaches only $M=3$; scaling costs $X^d$ \\

Integer Chebyshev & Directly controls integer coefficients & Values integral at all integer nodes & Original interval has capacity $\asymp X$; normalization loses integrality \\

$\Int(\Bsem,\Z)$ regular basis & Potentially large admissible denominators & Integral on every smooth node & No known regular basis has both sufficient denominator and Chebyshev norm \\

Dyadic interpolation & Exact rational certificate; short proofs & Fits $2^n\mapsto3^n$ exactly & Extrapolation denominator/error loses for large $M$ \\

$q$-Newton & Entire geometric interpolation; explicit coefficients & Integer-valued on a geometric axis & First extrapolation is a superexponentially large integer \\

Discrete minimax & Exact alternation and cofactor formula & Integral evaluation matrix & Clearing coefficients changes subunit error $D/S$ into integer error $D$ \\

Rational/Pad\'e & Can place poles near the branch set & Integer numerator/denominator at integer nodes & Locking scale is $1/|Q(q)|$; coefficients often lie in $\Q(\beta)$ \\

Hermite--Pad\'e exponential & Powerful determinant and multiplicity machinery & Works for algebraic differential data & Natural arguments and system coefficients are logarithmic/transcendental \\

Orthogonal / Christoffel & Exact Gram quotient; good conditioning & Determinants integral with integral weights & Real least squares ignores the evaluation-lattice coset \\

M\"untz system & Flexible semigroup of frequencies & Every term integral on $\Bsem$ under a counterexample & Integer coefficient space has a hard unit floor; signed exponents incur denominators \\
\bottomrule
\caption{Approximation-theoretic approaches and their audited failure points.}
\label{tab:comparison}
\end{longtable}
}

The common pattern is now transparent.  Real approximation can be made strong, and arithmetic
integrality can be made exact, but every known construction pays for moving from one to the
other.  The conjecture asks for a single auxiliary object that is simultaneously near-minimax
at infinity and denominator-rich at every finite prime.


\section{Conclusions}
\label{sec:conclusion}

No unconditional proof of the Alaoglu--Erd\H{o}s conjecture has been obtained in this
continuation.  The research completed here nevertheless narrows the polynomial-approximation
program to a small number of exact targets.

The sharp Archimedean result is Theorem~\ref{thm:optimal-block}: among all multiplicative
blocks, $c=3+2\sqrt2$ maximizes the product of smooth-node density and Bernstein-ellipse rate.
The optimized inequality can reach only the smallest hypothetical non-power generator
$M=3$; it is geometrically incapable of treating $M\ge5$.  This is a rigorous boundary on an
entire class of classical polynomial proofs.

The sharp arithmetic result is the unit barrier.  On smooth nodes, integer or conditional
M\"untz coefficients give integral residuals, so every unsuccessful approximation has norm at
least one.  The discrete cofactor formula shows exactly what happens to the real minimax
solution: its subunit error $D/S$ becomes the nonzero integer determinant $D$ when the
coefficient denominator is cleared.  The missing theorem is therefore not a slightly better
real approximation estimate.  It is an arithmetic construction that clears the approximant
without multiplying the target coefficient.

Two constructive by-products are exact.  Dyadic interpolation gives short symbolic
certificates excluding $M=3,5,6,7$, and the $q$-Newton basis gives a complete formula for
integer-valued interpolation on a geometric progression, including its superexponential
first-extrapolation error.  These results are modest compared with the inherited finite scan,
but they are independent, proof-producing, and formalization-ready.

The most promising next step is the generator-three branch in
$\Int(\Bsem,\Z)$.  Bhargava $p$-orderings and their recent extensions provide a language for
its allowable denominators; integer Chebyshev and potential theory provide the real norm to
beat.  A regular basis combining both at the optimal block would settle $3^{\thetaq}\notin\Z$.
To reach all $M$, a stronger ingredient is needed: a weighted rational construction, a signed
M\"untz basis with compensated denominators, or a genuinely new logarithmic auxiliary-function
theorem crossing the four-exponentials boundary.

\begin{resultbox}
\textbf{Final status.}
The conjecture remains open.  The continuation proves a sharp optimal-block theorem, exact
polynomial and rational locking principles, a cofactor minimax identity with an arithmetic
integerization theorem, a $q$-Newton interpolation theorem, and four symbolic generator
exclusions.  It also identifies the precise arithmetic approximation theorem that would turn
the only favorable Bernstein branch into a proof.
\end{resultbox}


\appendix

\section{Detailed calculus for the optimal block}
\label{app:optimal-block}

For reference, the ellipse parameter can be derived directly.  Map $[1,c]$ affinely to
$[-1,1]$ by
\[
 x=\frac{2t-(1+c)}{c-1}.
\]
The branch point $t=0$ maps to
\[
 x_0=-\frac{c+1}{c-1}<-1.
\]
The Bernstein ellipse through $x_0$ has parameter
\[
 \rho=|x_0|+\sqrt{x_0^2-1}
 =\frac{c+1+2\sqrt c}{c-1}
 =\frac{\sqrt c+1}{\sqrt c-1}.
\]
If $c=e^{2y}$, then
\[
 \rho=\frac{e^y+1}{e^y-1}=\coth(y/2)=e^u.
\]
The half-angle identity gives
\[
 \tanh(y/2)=e^{-u},
 \qquad
 \sinh y=\frac{2e^{-u}}{1-e^{-2u}}=\frac1{\sinh u}.
\]
Thus the optimization of $\Psi(c)=2yu$ is symmetric in $y$ and $u$ under the constraint
$\sinh y\sinh u=1$.  Strict convexity of $z\mapsto\log\sinh(e^z)$ gives
\[
 \log\sinh\sqrt{yu}
 \le\frac12\log\sinh y+\frac12\log\sinh u=0.
\]
Therefore $yu\le\arsinh(1)^2$, with equality only at $y=u=\arsinh(1)$.

The numerical constants are
\begin{align*}
 c_*&=5.8284271247461900976\ldots,\\
 \rho(c_*)&=2.4142135623730950488\ldots,\\
 \max\Psi&=1.5536387997913919630\ldots,\\
 M_{\mathrm{crit}}&=4.1131245860194105174\ldots.
\end{align*}


\section{Proof details for the generalized power zero bound}
\label{app:cheb-zero}

Let $\lambda_0<\cdots<\lambda_s$ be distinct real numbers.  The Wronskian of
$t^{\lambda_0},\ldots,t^{\lambda_s}$ on $(0,\infty)$ is
\[
 W(t)=t^{\lambda_0+\cdots+\lambda_s-s(s+1)/2}
      \prod_{0\le i<j\le s}(\lambda_j-\lambda_i).
\]
It never vanishes.  Hence the functions form an extended complete Chebyshev system, and any
nontrivial linear combination has at most $s$ positive zeros counted with multiplicity.
This proves every zero-count statement used in the polynomial, rational, and M\"untz locking
theorems.

An elementary induction is also available.  Divide a nonzero generalized polynomial by its
smallest power, differentiate, and apply Rolle's theorem.  Each differentiation removes the
constant term and leaves one fewer exponent.  This proves the distinct-zero form without the
full Wronskian formalism.


\section{Derivation of the discrete cofactor formula}
\label{app:cofactor}

Let $A$ be an $(N+1)\times N$ real matrix of rank $N$, $f\in\R^{N+1}$, and
$V=[A\ f]$.  For each $i$, let
\[
 C_i=(-1)^{i+N}\det A_{\widehat i},
\]
where $A_{\widehat i}$ deletes row $i$.  Laplace expansion gives
\[
 A^TC=0,
 \qquad
 C^Tf=\det V=D.
\]
The dual of the distance from $f$ to $\operatorname{im}A$ in $\ell^\infty$ is
\[
 \max\{y^Tf:A^Ty=0,\ \|y\|_1\le1\}.
\]
The nullspace is spanned by $C$, so the value is $|D|/S$ with
$S=\sum_i|C_i|$.  If the ordered evaluation matrix is strictly totally positive, then the
$C_i$ alternate in sign.  Setting $s_i=\sgn(C_i)$, the square matrix $B=[A\ s]$ has
\[
 \det B=C^Ts=S.
\]
The equation
\[
 Ac+\frac DSs=f
\]
therefore has a unique solution and is the alternant.  If $A,f$ are integral, multiply by
$S$:
\[
 B\binom{Sc}{D}=Sf.
\]
Since $\operatorname{adj}(B)$ is integral,
\[
 \binom{Sc}{D}=\operatorname{adj}(B)f\in\Z^{N+1}.
\]
This proves Theorems~\ref{thm:cofactor-minimax} and~\ref{thm:arithmetic-alternant}.


\section{The geometric Newton identity}
\label{app:qnewton}

For $q\ge2$ and $n\ge k$, one has
\[
 B_k^{(q)}(q^n)
 =\prod_{i=0}^{k-1}\frac{q^n-q^i}{q^k-q^i}
 =\prod_{i=0}^{k-1}\frac{q^{n-i}-1}{q^{k-i}-1}
 =\qbinom{n}{k}{q}.
\]
For $n<k$, the product vanishes because one factor is $q^n-q^n$.
The Newton coefficient at stage $k$ is the error at the new node $q^k$, hence
\[
 a_k=r^k-P_{k-1}(q^k).
\]
Induction using the Gaussian-binomial recursion gives
\[
 a_k=\prod_{i=0}^{k-1}(r-q^i).
\]
Equivalently, the exact identity
\[
 r^n=\sum_{k=0}^n\qbinom{n}{k}{q}
             \prod_{i=0}^{k-1}(r-q^i)
\]
is the Newton expansion of the data $q^n\mapsto r^n$.  Truncating at $d$ gives
Theorem~\ref{thm:q-newton}; at $n=d+1$ only one tail term survives.


\section{Literature chronology and relevance matrix}
\label{app:literature}

\begin{longtable}{p{0.16\textwidth}p{0.28\textwidth}p{0.47\textwidth}}
\toprule
Date & Source & Relevance to this continuation \\
\midrule
\endfirsthead
\toprule
Date & Source & Relevance to this continuation \\
\midrule
\endhead
1944 & Alaoglu--Erd\H{o}s & Original arithmetic question in the theory of highly composite numbers. \\
1957 & Schneider & Places the surrounding four-exponentials problem among central open transcendence questions. \\
1997 & Bhargava & $p$-orderings and generalized factorials for integer-valued polynomial rings on arbitrary subsets. \\
2003 & Borwein--Pinner--Pritsker & Monic integer Chebyshev problem and exact constants on arithmetic compact sets. \\
2005 & Pritsker & Integer Chebyshev polynomials, potential theory, capacity, and small integer-coefficient polynomials. \\
2023 & Waldschmidt & Modern survey explicitly connecting the question to four exponentials and Schanuel-type conjectures. \\
2024/2026 & Dasgupta--Kakde & Matrix Coefficient Conjecture and its relation to strong four exponentials; published in Advances in Mathematics. \\
2025 & Massold & Transcendence-degree estimates for exponentials of products; stronger cases reduced to torus intersections. \\
2025 & Fischler--Rivoal & New exponential-value transcendence measure using Hermite--Pad\'e approximants and an optimized Siegel determinant. \\
2025 & Adamczewski--Faverjon & Explicit algebraic-independence measures for values of $E$- and $M$-functions at algebraic points. \\
2025/2026 & Gorbachev--Ivanov--Tikhonov & Discrete Chebyshev alternation and Sturm oscillation for discrete polynomial systems. \\
2025/2026 & Berger--Sprang & Links regular bases of integer-valued polynomial rings with $p$-adic Fourier theory. \\
2025 & Lagarias--Yangjit & Extends Bhargava-style orderings and generalized factorials from prime ideals to general ideals. \\
\bottomrule
\caption{Selected literature used to delimit the approximation program.}
\label{tab:literature}
\end{longtable}

The chronology shows that the relevant tools are progressing on both sides: quantitative
auxiliary functions at the Archimedean place and generalized integer-valued polynomial bases
at finite places.  What remains absent is a theorem joining them for the specialized
$2\times2$ logarithmic matrix.


\begin{thebibliography}{99}

\bibitem{AE1944}
L.~Alaoglu and P.~Erd\H{o}s,
\newblock \emph{On highly composite and similar numbers},
\newblock Transactions of the American Mathematical Society \textbf{56} (1944), 448--469.
\newblock DOI: \href{https://doi.org/10.1090/S0002-9947-1944-0011087-2}{10.1090/S0002-9947-1944-0011087-2}.

\bibitem{Schneider1957}
Th.~Schneider,
\newblock \emph{Einf\"uhrung in die transzendenten Zahlen},
\newblock Springer, 1957.

\bibitem{Waldschmidt2000}
M.~Waldschmidt,
\newblock \emph{Diophantine Approximation on Linear Algebraic Groups},
\newblock Grundlehren der mathematischen Wissenschaften 326, Springer, 2000.

\bibitem{Waldschmidt2023}
M.~Waldschmidt,
\newblock \emph{The Four Exponentials Problem and Schanuel's Conjecture},
\newblock in J.-M.~Morel and B.~Teissier (eds.), \emph{Mathematics Going Forward},
Lecture Notes in Mathematics 2313, Springer, 2023, pp.~579--592.
\newblock DOI: \href{https://doi.org/10.1007/978-3-031-12244-6_39}{10.1007/978-3-031-12244-6\_39}.

\bibitem{DasguptaKakde2024}
S.~Dasgupta and M.~Kakde,
\newblock \emph{Ranks of Matrices of Logarithms of Algebraic Numbers II: The Matrix
Coefficient Conjecture},
\newblock arXiv:2408.08178, 2024; subsequently published in 2026 as cited below.

\bibitem{DasguptaKakde2026}
S.~Dasgupta and M.~Kakde,
\newblock \emph{Ranks of matrices of logarithms of algebraic numbers II: The Matrix
Coefficient Conjecture},
\newblock Advances in Mathematics \textbf{487} (2026), Article 110753.
\newblock DOI: \href{https://doi.org/10.1016/j.aim.2025.110753}{10.1016/j.aim.2025.110753}.

\bibitem{Massold2025}
H.~Massold,
\newblock \emph{Transcendence degrees of fields generated by exponentials of products},
\newblock arXiv:2506.01123, 2025.

\bibitem{FischlerRivoal2025}
S.~Fischler and T.~Rivoal,
\newblock \emph{A new transcendence measure for the values of the exponential function at algebraic arguments},
\newblock arXiv:2502.17992v1, February 2025.

\bibitem{AdamczewskiFaverjon2025}
C.~Faverjon and B.~Adamczewski,
\newblock \emph{Algebraic Independence Measures for Values of $E$-functions and
$M$-functions},
\newblock arXiv:2502.09999, 2025.

\bibitem{GorbachevIvanovTikhonov2025}
D.~V.~Gorbachev, V.~I.~Ivanov, and S.~Yu.~Tikhonov,
\newblock \emph{Chebyshev systems and Sturm oscillation theory for discrete polynomials},
\newblock Forum of Mathematics, Sigma \textbf{14} (2026), e84.
\newblock DOI: \href{https://doi.org/10.1017/fms.2026.10168}{10.1017/fms.2026.10168};
\newblock arXiv:2501.02358v2, August 2026.

\bibitem{Bhargava1997}
M.~Bhargava,
\newblock \emph{$P$-orderings and polynomial functions on arbitrary subsets of Dedekind
rings},
\newblock Journal f\"ur die reine und angewandte Mathematik \textbf{490} (1997), 101--127.

\bibitem{LagariasYangjit2025}
J.~C.~Lagarias and W.~Yangjit,
\newblock \emph{$B$-orderings for all ideals $B$ of Dedekind domains and generalized
factorials},
\newblock arXiv:2502.19072, 2025.

\bibitem{BergerSprang2025}
L.~Berger and J.~Sprang,
\newblock \emph{Integer-valued polynomials and $p$-adic Fourier theory},
\newblock Rendiconti del Seminario Matematico della Universit\`a di Padova,
Online First, July 21, 2026.
\newblock DOI: \href{https://doi.org/10.4171/RSMUP/204}{10.4171/RSMUP/204};
\newblock arXiv:2502.18053v2, April 2025.

\bibitem{Pritsker2005}
I.~E.~Pritsker,
\newblock \emph{Small polynomials with integer coefficients},
\newblock Journal d'Analyse Math\'ematique \textbf{96} (2005), 151--190.
\newblock DOI: \href{https://doi.org/10.1007/BF02787827}{10.1007/BF02787827}.

\bibitem{BorweinPinnerPritsker2003}
P.~B.~Borwein, C.~G.~Pinner, and I.~E.~Pritsker,
\newblock \emph{Monic integer Chebyshev problem},
\newblock Mathematics of Computation \textbf{72} (2003), no.~244, 1901--1916.
\newblock DOI: \href{https://doi.org/10.1090/S0025-5718-03-01477-7}{10.1090/S0025-5718-03-01477-7}.

\bibitem{KarlinStudden1966}
S.~Karlin and W.~J.~Studden,
\newblock \emph{Tchebycheff Systems: With Applications in Analysis and Statistics},
\newblock Interscience, 1966.

\bibitem{Lorentz1986}
G.~G.~Lorentz,
\newblock \emph{Approximation of Functions},
\newblock second edition, Chelsea, 1986.

\bibitem{Trefethen2013}
L.~N.~Trefethen,
\newblock \emph{Approximation Theory and Approximation Practice},
\newblock SIAM, 2013.

\bibitem{CahenChabert1997}
P.-J.~Cahen and J.-L.~Chabert,
\newblock \emph{Integer-Valued Polynomials},
\newblock Mathematical Surveys and Monographs 48, American Mathematical Society, 1997.

\bibitem{UnifiedReport2026}
The ProveIt project,
\newblock \emph{The Alaoglu--Erd\H{o}s Integer-Exponent Conjecture: Machine-Checked Partial
Results, Exact Conditional Structure, Determinant Methods, and the Remaining Barrier},
\newblock supplied unified report, August 21, 2026.

\end{thebibliography}

\end{document}
